%% Apuntes de Fotodeteccion y Trayectorias Cuanticas
%
%\documentclass[twocolumn]{article}
\documentclass[twocolumn,aps,eqsecnum,showkeys,amsmath]{revtex4}
\usepackage[spanish,activeacute]{babel}
\usepackage{amsmath}
\usepackage{graphicx}       
\usepackage{color}
\usepackage{cancel}
\spanishdecimal{.}
\usepackage[colorlinks,urlcolor=blue,citecolor=blue,linkcolor=blue]{hyperref}% add hypertext capabilities


% vectors (the most common ones)
\newcommand{\mbi}{\mathbf{i}}
\newcommand{\mbj}{\mathbf{j}}
\newcommand{\mbk}{\mathbf{k}}
\newcommand{\mbp}{\mathbf{p}} 
\newcommand{\mbq}{\mathbf{q}}
\newcommand{\mbr}{\mathbf{r}}
%\newcommand{\mbrp}{\mathbf{r}^{\prime}}
\newcommand{\mbv}{\mathbf{v}}
\newcommand{\mbE}{\mathbf{E}}
\newcommand{\mbR}{\mathbf{R}}

\newcommand{\mcC}{\mathcal{C}} 
\newcommand{\mcH}{\mathcal{H}} 	%Hamiltonian
\newcommand{\mcL}{\mathcal{L}} 	
\newcommand{\mcS}{\mathcal{S}} 
\newcommand{\mcE}{\mathcal{E}} 
\newcommand{\mcV}{\mathcal{V}} 	%Hamiltonian
\newcommand{\Tr}{\mathrm{Tr}}
\newcommand{\qed}{\mathrm{Q.E.D.}}
\newcommand{\p}{\partial}
%\newcommand{\lan}{\langle}     \newcommand{\ran}{\rangle}
\newcommand{\adg}{a^{\dag}}
\newcommand{\Adg}{A^{\dag}} 	
\newcommand{\Hdg}{H^{\dag}} 	 
\newcommand{\Udg}{U^{\dag}} 	

\newcommand{\lan}{\langle}
\newcommand{\ran}{\rangle}
%\newcommand{\adg}{a^{\dag}}

\providecommand{\abs}[1]{\left\lvert#1\right\rvert}
\providecommand{\norm}[1]{\lVert #1 \rVert}
\providecommand{\vxp}[1]{\langle #1 \rangle}
\providecommand{\bra}[1]{\langle #1 \rvert}
\providecommand{\ket}[1]{\lvert #1 \rangle}
\providecommand{\bbra}[1]{\langle #1 \rvert\rvert}
\providecommand{\kket}[1]{\lvert\lvert #1 \rangle}
\providecommand{\kkket}[1]{\lvert #1 \rangle\rangle}
\providecommand{\braket}[2]{\langle #1 \rvert #2 \rangle}
%\providecommand{\vxp3}[3]{\langle #1 |#2| #3 \rangle}
\providecommand{\tr}[1]{\text{tr}\left[ #1 \right]}
\newcommand{\ketbra}[2]{\left| {#1} \right\rangle\left\langle {#2}\right|}
%\newcommand{\mmax}{_{\text{max}}}

\newcommand{\mru}[1]{\mathrm{#1}} %unidades fisicas
\newcommand{\inc}[1]{\mathrm{(#1)}} %incisos en formulas desplegadas
\newcommand{\tbn}[1]{\textbf{#1}} %numero de problema

\begin{document}

\title{Apuntes para el curso  \\
Teor\'ia de Fotodetecci\'on y Trayectorias Cu\'anticas}
\author{H\'ector M. Castro Beltr\'an}
\affiliation{Centro de Investigaci\'on en Ingenier\'ia y 
Ciencias Aplicadas, IICBA, \\ Universidad Aut\'onoma del Estado de Morelos,  
Av. Universidad 1001, 62209 Cuernavaca, Morelos, M\'exico}

\date{\today}
\begin{abstract} 
Este curso trata de la simulaci\'on de mediciones de luz mediante 
ecuaciones estoc\'asticas de Schr\"odinger, ecuaciones para la evoluci\'on 
de la funci\'on de onda. Estos m\'etodos son llamados Trayectorias 
Cu\'anticas o de Ecuaci\'on de Onda de Monte Carlo, y engloban m\'etodos 
de saltos cu\'anticos y difusi\'on de estados cu\'anticos para simular 
fotoconteo, detecci\'on homodina y heterodina y  detecci\'on de espectros. 
Trataremos solamente aspectos b\'asicos de la teor\'ia de la fotodetecci\'on 
de inter\'es inmediato para las simulaciones. La mayor\'ia de los ejemplos 
usan como base la teor\'ia de fluorescencia resonante de un \'atomo de 
dos niveles. 

Apuntes en proceso, tomados de varias fuentes. La notaci\'on entre 
secciones no es uniforme, hay material repetido, mal ubicado, partes en 
ingl\'es y otras en espa\~nol, etc. En estas notas no intentamos una 
redacci\'on demasiado precisa; nos enfocamos en mencionar los temas 
y puntos m\'as importantes. 

Un apunte sobre teor\'ia de coherencia y espectro estacionario 
complementa el curso. 
\end{abstract}

\keywords{ Photon-counting, homodyne and heterodyne detection, 
quantum trajectories, quantum jumps, quantum state diffusion, 
waiting-time distribution, intensity correlations, amplitude-intensity 
correlation, resonance fluorescence,}
\maketitle

%%%%%%%%10%%%%%%%%20%%%%%%%%30%%%%%%%%40
%%%%%%%%50%%%%%%%%60%%%%%%%%70%%%
\section{Introducci\'on}
Informaci\'on importante sobre la estad\'istica de un campo se obtiene 
simplemente contando los fotopulsos producidos en el detector. El 
n\'umero de pulsos var\'ia si repetimos el experimento una y otra vez 
porque el proceso de emisi\'on fotoel\'ectrica es fundamentalmente 
probabil\'istico. Tres factores contribuyen a este hecho: i) las 
fluctuaciones estad\'isticas del campo detectado, ii) el car\'acter 
cu\'antico de la interacci\'on entre detector y campo, y 
iii) imperfecciones en el proceso de medición. 

La tarea de la teor\'ia de detecci\'on fotoel\'ectrica (DF) es establecer 
una relaci\'on entre un haz de luz y una secuencia de emisi\'on de 
fotoelectrones. En la teor\'ia DF semicl\'asica la emisi\'on de 
fotoelectrones esta gobernada por la intensidad estoc\'astica cl\'asica 
del campo $I(t)$. El detector produce una corriente an\'aloga que 
describe un segundo proceso estoc\'astico $i(t)$. La teor\'ia de DF 
debe relacionar $i(t)$ con $I(t)$. 

%%%%%%%%10%%%%%%%%20%%%%%%%%30%%%%%%%%40
%%%%%%%%50%%%%%%%%60%%%%%%%%70%%%
\section{Fotoconteo y Fotocorriente}

%%%%%%%%10%%%%%%%%20%%%%%%%%30%%%%%%%%40
\subsection{Distribuci\'on de Conteo de Fotoelectrones} %%%%%%%%
La DCF es la densidad de probabilidad $p(n,t,T)$ de contar $n$ 
fotoelectrones en el intervalo $(t,t+T]$. Calcularemos primero la DCF 
para un campo incidente cl\'asico de intensidad constante $\bar{I}$, es 
decir, sin fluctuaciones, y tiene frecuencia $\omega$. Este ejemplo 
permite separar los efectos de las estad\'isticas del campo y del detector. 
El haz es intersectado por un detector de \'area $A$. Sea 
$A \bar{I}/\hbar \omega$ el n\'umero promedio de fotones que entran 
al detector por segundo. Sea entonces $A \bar{I} T/\hbar \omega$ el 
n\'umero promedio de fotones contados en el tiempo $T$, suponiendo 
que cada fot\'on es convertido en un fotoelectr\'on. Esta es una 
idealizaci\'on. El detector tiene una eficiencia cu\'antica $\eta$ 
($0< \eta <1$), que es la proporci\'on de fotones convertidos en 
fotoelectrones. Entonces, 
%
\begin{eqnarray}
\eta \frac{A \bar{I}}{\hbar \omega} T =\xi \bar{I} T \,, 
\end{eqnarray}
%%
es el n\'umero promedio de fotones contados en el tiempo $T$, donde 
$\xi =\eta A/\hbar \omega$.  

El tratamiento cu\'antico del efecto fotoel\'ectrico proporciona la 
probabilidad de emisi\'on fotoel\'ectrica en un tiempo de interacci\'on 
infinitesimal $\Delta t$. Dividamos $T$ en $N=T/\Delta t \gg 1$ 
subintervalos. Si el campo es constante todos los subintervalos 
$\Delta t$ son iguales (equivalentes), suponiendo que el detector 
se recupera inmediatamente (ignorando el tiempo muerto) despu\'es 
de emitir un fotoelectr\'on, por lo que queda disponible para emitir otro 
inmediatamente. 

Obtendremos la DCF de dos maneras.

\subsubsection{Distribuci\'on  de Probabilidad de Fotoconteo 1 
\label{sec:DPF1} } %%%
Calculemos la probabilidad de contar $n$ fotoelectrones en el tiempo 
$T$ de la estad\'istica de $N$ volados independientes, \textit{\'aguila} 
si se emite un electr\'on en $\Delta t$, \textit{sol} si no se emite. 
La probabilidad $p$ para \'aguilas es proporcional al flujo de fotones 
incidentes sobre el detector, $p=\xi \bar{I} \Delta t$; la 
probabilidad para soles es $q=1-p=1-\xi \bar{I} \Delta t$. Al hacer 
$\Delta t$ suficientemente peque\~no eliminamos eventos con dos o mas 
fotoelectrones emitidos en $\Delta t$. La probabilidad para $n$ 
\'aguilas y $N-n$ soles esta dada por la distribuci\'on binomial 
%
\begin{eqnarray*}
p(n,t,T) &=& \frac{N!}{n!(N-n)!} p^n q^{N-n}     \\ 
	&=& \frac{N(N-1) \cdots(N-n+1)}{n!} p^n (1-p)^{N-n} \\
&=& \left(1-\frac{1}{N} \right) \cdots \left(1-\frac{n-1}{N} \right) \\ 
	&& \times \frac{ (\xi \bar{I} \Delta t N)^n }{n!}
	(1- \xi \bar{I} \Delta t)^{N-n} \,.
\end{eqnarray*}
%%
Tomando los l\'imites $N \to \infty$ y $\Delta t \to 0$ con 
$N \Delta t =T =$ cte 
%
\begin{eqnarray*}
(1- \frac{\xi \bar{I} T}{N})^{N-n} 
&=& 1-(N-n)\frac{\xi \bar{I} T}{N} \\ 
  && +\frac{(N-n)(N-n-1)}{2!} \left(\frac{\xi \bar{I} T}{N} \right)^2 
	-\cdots \\
&\simeq& 1- \xi \bar{I} T +\frac{(\xi \bar{I} T)^2}{2!} -\cdots 
	=e^{-\xi \bar{I} T } \,.
\end{eqnarray*}
%%
Entonces, 
%
\begin{eqnarray}
p(n,t,T) &=& p(n,T) 
	=\frac{(\xi \bar{I} T)^n}{n!} e^{-\xi \bar{I} T } \,,
\end{eqnarray}
%%
que es la distribuci\'on de Poisson, $P_n =\bar{n}^n e^{-\bar{n}}/n!$, 
donde $\bar{n}= \xi \bar{I} T$ es el n\'umero promedio de 
fotoelectrones emitidos en el tiempo $T$. Esta es independiente 
del inicio del fotoconteo, dado que con un campo constante todos 
los intervalos $T$ son iguales. Las fluctuaciones Poissonianas en 
el fotoconteo bajo iluminaci\'on constante son el origen del 
denominado \textit{ruido de disparo}.
 
Notas: i) El ruido del campo es denominado \textit{ruido de fotones}. 
Si en un sistema de detecci\'on el ruido de fotones es dominante, 
el sistema es considerado \textit{ideal}, en el sentido de que es 
imposible mejorar su funcionamiento. ii) El ruido de fondo de un campo 
puede ser muy importante, en particular, ruido de radiaci\'on t\'ermica 
(p. ej., infrarrojo por temperatura del instrumento \'optico) o el ruido 
de fondo c\'osmico (temperatura $T \sim 2.7$K). iii) Hay ruido por 
fluctuaciones en la \textit{corriente obscura} de un fotomultiplicador. 
iv) Hay ruido de \textit{Johnson}, producido por fluctuaciones de 
voltaje en resistencias y es de origen t\'ermico.

Si dejamos que la intensidad cambie en el tiempo, contando fotones 
durante un tiempo $T$, la f\'ormula de fotoconteo a\'un es v\'alida si 
hacemos 
%
\begin{subequations}
\begin{eqnarray}
\bar{I} &=& \frac{1}{T} \int_0^T I(t) dt \,, \\
P_n(t) &=& \frac{(\xi \bar{I} T)^n}{n!} \,e^{-\xi \bar{I} T}  	\,, 
	\label{eq:MamdelPhotocount}
\end{eqnarray}
\end{subequations}
%%
conocida como la \textit{F\'ormula de Mandel de fotoconteo}. 

\subsubsection{Distribuci\'on  de Probabilidad de Fotoconteo 2} %%%%
Tomado de AWH, Rice (p.7.2) y Fox (p.90).  Pendiente: uniformizar  
notaci\'on con lo anterior. 
\newline Supongamos que tenemos un campo EM cl\'asico de intensidad 
$I$ sin fluctuaciones y que tenemos un un detector que genera una 
se\~nal proporcional a esa intensidad. Digamos que el detector tiene una 
probabilidad de registrar un evento de fotodetecci\'on (hacer click) sobre 
un intervalo de tiempo $dt$ dada por $\eta I dt$. Al tiempo $t+dt$, si 
hemos detectado $n$ fotones, con probabilidad $P_n(t+dt)$, esto pudo 
ocurrir de dos maneras: podr\'iamos haber detectado $n-1$ fotones en el 
intervalo $[0,t]$, y un fot\'on en el intervalo $[t,t+dt]$, o que detectamos 
$n$ fotones en el intervalo $[0,t]$ y ninguno en el siguiente intervalo 
$[t,t+dt]$. Esto se expresa como 
%
\begin{eqnarray*}
P_n(t+dt) &=& P_{n-1}(t) \eta I dt +P_n(t) (1-\eta I dt) \,.
\end{eqnarray*}
%%
Reacomodando t\'erminos, con $\mu= \eta I$, 
%
\begin{eqnarray*}
\frac{d P_n(t)}{dt} &=& \frac{ P_n(t+dt) -P_n(t)}{dt} 
	=\mu P_{n-1}(t) -\mu P_n(t)  \,.
\end{eqnarray*}
%%
Para resolver esta ecuaci\'on tomamos la condici\'on inicial $P_0(0)=1$ 
y $P_{n\geq 1}(0) =0$. Para $n=0$, 
%
\begin{eqnarray*}
\frac{d P_0(t)}{dt} &=& -\mu P_0(t) 		\\
\frac{d P_0(t)}{P_0(t)} &=& -\mu dt \,, \qquad \to P_0(t)= e^{-\mu t} \,.
\end{eqnarray*}
%%
Para $n=1$ [con $P_1(0)=0$], 
%
\begin{eqnarray*}
\frac{d P_1(t)}{dt} &+&\mu P_1(t) = \mu P_0(t) =\mu \,e^{-\mu t} \,,	\\	
\frac{d}{dt} \left[ e^{\mu t} P_1  \right] &=& \mu 	
	\quad \to e^{\mu t} P_1(t) = \mu t +c_1 	\,, \\ 
P_1(t) &=& \mu t \,e^{-\mu t} 	\,.
\end{eqnarray*}
%%

Para $n=2$ [con $P_2(0)=0$], 
%
\begin{eqnarray*}
\frac{d P_2(t)}{dt} &+&\mu P_2(t) = \mu P_1(t) =\mu^2 t \,e^{-\mu t}  \,, \\	
\frac{d}{dt} \left[ e^{\mu t} P_2  \right] &=& \mu^2 t 	\,, 
	\quad \to e^{\mu t} P_2(t) = \frac{\mu^2 t^2}{2} +c_2 	\,,\\ 
P_2(t) &=& \frac{\mu^2 t^2}{2} \,e^{-\mu t} 	\,.
\end{eqnarray*}
%%

Para $n=3$ [con $P_3(0)=0$], 
%
\begin{eqnarray*}
\frac{d P_3(t)}{dt} &+&\mu P_3(t) = \mu P_2(t) 
	=\frac{\mu^3 t^2}{2} \,e^{-\mu t}  	\,, \\	
\frac{d}{dt} \left[ e^{\mu t} P_3  \right] &=& \frac{\mu^3 t^2}{2} 	\,,
	\quad \to e^{\mu t} P_3(t) = \frac{\mu^3 t^3}{2 \cdot 3} +c_3 \,,\\ 
P_3(t) &=& \frac{\mu^3 t^3}{2 \cdot 3} e^{-\mu t} 	\,,	\\
&\vdots& 	\\
P_n(t) &=& \frac{\mu^n t^n}{n!} \,e^{-\mu t}  	\,.
\end{eqnarray*}
%%
Esta es la distribuci\'on de Poisson con n\'umero promedio de 
fotoconteo $\mu = \eta I$ y dispersi\'on $\sqrt{\mu}$. 

%%%%%%%%%%%%%%%%%%%%%%%%%%%%%%%%
\subsection{Fotocorriente en Detecci\'on Homodina} 
\textcolor{red}{Resumen de la secci\'on 9.1 de \cite{Carmichael93} 
en preparaci\'on.}

En esta subsecci\'on relajaremos la condici\'on de que un intervalo 
$\Delta t$ sea tan corto que solo un fotopulso puede ser emitido. Ahora 
$\Delta t$ es suficientemente largo para que muchos fotoelectrones 
sean emitidos, pero muy corto comparado con la escala de fluctuaciones 
del campo de la fuente. Si $\xi \bar{I} \Delta t$ es un n\'umero grande las 
fluctuaciones est\'an caracterizadas por la distribuci\'on de Gauss,    
%
\begin{eqnarray}
p(n_t,t,t+\Delta t) &=& \frac{1}{\sqrt{ 2\pi \xi \bar{I}(t) \Delta t}} 	\nonumber \\
	&& \times \exp{ \left( - \frac{(n_t -\xi \bar{I}(t) \Delta t)^2}
	{2\xi \bar{I}(t) \Delta t} \right) } \,. \qquad
\end{eqnarray}
%%
Entonces, la carga $\Delta Q$ emitida por el fotoc\'atodo en el intervalo 
$\Delta t$ esta dada por 
 %
\begin{eqnarray}
\frac{\Delta Q}{e} = \xi \bar{I}(t) \Delta t +\sqrt{ \xi \bar{I}(t)} \Delta W \,, 
\end{eqnarray}
%%
donde $\Delta W$ es un incremento de Wiener. La fotocorriente es ahora 
 %
\begin{eqnarray} 	\label{eq:photocurrent}
\frac{i(t)}{Ge} = \frac{\Delta Q}{G e \Delta t} 
	= \xi \bar{I}(t) +\sqrt{ \xi \bar{I}(t)} \eta_W(t) \,, 
\end{eqnarray}
%%
donde $G$ es un factor de ganancia y $\eta_W(t)$ indica ruido blanco 
Gaussiano: 
%
\begin{equation} 	
\overline{dW} =0 \,,  \qquad \overline{\eta_W(t) \eta_W(t')} = \delta(t-t') \,.
\end{equation}
%%
El t\'ermino $\sqrt{ \xi \bar{I}(t)} \eta_W(t)$ es el ruido de disparo. 
Vemos entonces que la fotocorriente $i(t)$ no es una simple r\'eplica 
de la intensidad \'optica.

Veamos ahora la detecci\'on homodina (DHm). En la DHm el flujo de 
fotones visto por el detector es la superposici\'on de la se\~nal 
\'optica de amplitud fluctuante $e_s(t)$ y la amplitud constante 
$E_{lo} = |E_{lo}| e^{i\phi}$ del oscilador local (OL),  
%
\begin{equation} 	
\xi \bar{I}(t) = \xi |E_{lo} +e_s(t)|^2 \,.  
\end{equation}
%%
Sustituyendo $\xi \bar{I}(t)$ en la Ec.(\ref{eq:photocurrent}) y con la 
suposici\'on de que $|E_{lo}| \gg |e_s(t)|$ tenemos que 
 %
\begin{eqnarray}
\frac{i(t)}{Ge} &\simeq& \xi |E_{lo}|^2 +\sqrt{\xi} \,|E_{lo}| \eta_W(t) 
	\nonumber \\
	&& +\xi |E_{lo} e_s^\ast(t) + E_{lo}^\ast e_s(t)| \,,
\end{eqnarray}
%%
donde solo conservamos t\'erminos de ruido al orden dominante en la 
amplitud del OL. $\eta_W$ es ruido blanco Gaussiano de la emisi\'on 
aleatoria de fotoelectrones a raz\'on promedio dominada por el flujo 
del OL. La se\~nal $e_s(t)$ tambi\'en introduce ruido, cuyo origen es 
la fuente que produce $e_s(t)$. Desde el punto de vista de la teor\'ia 
semicl\'asica de fotodetecci\'on, ambas fuentes de ruido son 
estad\'isticamente independientes. Entonces las fluctuaciones de la 
fotocorriente $\Delta (t)/Ge = i(t)/Ge -\xi |E_{lo}|^2$ estan caracterizadas 
por la funci\'on de correlaci\'on
 %
\begin{eqnarray} \label{eq:PhotoCurrentCorrel}
\frac{ \overline{\Delta i(t) \Delta i(t+\tau)} }{Ge} 
	&=& \xi |E_{lo}|^2 \overline{\eta_W(t) \eta_W(+\tau')} 	  \nonumber \\
	&& +4\xi^2  |E_{lo}|^2 \overline{ e_s^{\phi}(t) e_s^{\phi}(t+\tau) } 
	\nonumber \\
&=& \xi |E_{lo}|^2 \delta(\tau) 	\nonumber \\ 
	&&+4\xi^2  |E_{lo}|^2 \overline{ e_s^{\phi}(t) e_s^{\phi}(t+\tau) } \,, \qquad 
\end{eqnarray}
%%
donde 
%
\begin{equation} 	
e_s^{\phi}(t)  = \frac{1}{2} \left(  e_s(t) e^{-i\phi} + e_s^\ast(t) e^{i\phi} \right) \,, 
\end{equation}
%%
y $\phi$ es la fase del campo del oscilador local. La transformada 
de Fourier de la Ec.(\ref{eq:PhotoCurrentCorrel}) da el espectro de 
fluctuaciones de la fotocorriente. El primer t\'ermino es el espectro plano 
de ruido de disparo. El segundo t\'ermino \textit{a\~nade} ruido arriba al 
ruido de disparo. Esta es una deficiencia de la teor\'ia semicl\'asica ya 
que la luz comprimida hace que el ruido est\'e debajo del ruido de 
disparo. La soluci\'on, que veremos m\'as adelante, es que el ruido de la 
se\~nal debe tratarse de manera autoconsistente con el del detector; 
ambos ruidos estan correlacionados. Para ello, la teor\'ia debe ser 
completamente mec\'anico-cu\'antica.  

%%%%%%%%10%%%%%%%%20%%%%%%%%30%%%%%%%%40
%%%%%%%%50%%%%%%%%60%%%%%%%%70%%%
\section{Probabilidad de Detecci\'on}
%%%%%%%%%%%%%%%%%%%%%%%%%%%%%%%%
\subsection{Probabilidad de Detecci\'on de un Fot\'on} 
Consideremos un detector en equilibrio t\'ermico a temperatura ambiente 
$T \sim 300$ K. La probabilidad de que los \'atomos sean excitados por la 
energ\'ia t\'ermica es muy peque\~na [de la distribuci\'on de Boltzmann 
$P_B \propto \exp{(-\hbar \omega/k_B T)} \sim \exp{(-40)} \to 0$]. Por lo 
tanto, los \'atomos inician en su estado inferior y solo absorben radiaci\'on, 
es decir, solo el operador de destrucci\'on juega un papel: 
%
\begin{eqnarray}
\hat{E}^+ (x) \ket{i} \propto a \ket{i} \propto \ket{f} \,, 
\end{eqnarray}
%%
donde $\ket{i}$ y $\ket{f}$ son el estado inicial y final, respectivamente, 
del detector. La probabilidad de transici\'on por absorci\'on de un fot\'on 
ubicado en  $\mbr$ al tiempo $t$, en $x= \mbr, t$, es proporcional a 
%
\begin{eqnarray}
W = W_{i \to f} &=& |\bra{f} \hat{E}^+ (x) \ket{i} |^2 	\nonumber \\
	&=&  \bra{i} \hat{E}^- (x) \ket{f} \bra{f} \hat{E}^+ (x) \ket{i} \,, 
\end{eqnarray}
%%
Aunque sabemos que el \'atomo inicia en el estado base, normalmente 
no conocemos el estado inicial del campo, incluyendo variaciones 
estad\'isticas sobre $\ket{i}$. No nos interesa el estado final, de hecho, 
puede haber muchos de ellos, por lo que escribimos  
%
\begin{eqnarray}
W &=& \sum_f |\bra{f} \hat{E}^+ (x) \ket{i} |^2 
  = \sum_f \bra{i} \hat{E}^- (x) \ket{f} \bra{f} \hat{E}^+ (x) \ket{i} \nonumber \\ 
&=& \bra{i} \hat{E}^- (x) \left[ \sum_f \ket{f} \bra{f} \right] 
    \hat{E}^+ (x) \ket{i} 	\nonumber \\ 
    &=& \bra{i} \hat{E}^- (x)  \hat{E}^+ (x) \ket{i} . 
\end{eqnarray}
%%

 Usando el operador de densidad 
$\rho = \sum_i P_i \ketbra{i}{i}$ tenemos 
$W= \tr{\rho \hat{E}^- (x)  \hat{E}^+ (x)}$, que es la funci\'on de correlaci\'on 
de primer orden!

\subsection{Probabilidad de Detecci\'on de Dos Fotones} %%%%%%%%
Consideremos un \'atomo de dos niveles $|g\rangle$ y $|e\rangle$ 
excitado por un l\'aser monocrom\'atico. La emisi\'on de un fot\'on por 
emisi\'on espont\'anea viene asociada con la proyecci\'on del \'atomo al 
estado base: $A_{ge} |\Psi \rangle =|g\rangle$, donde 
$A_{ge} = |g \rangle \langle e|$. El experimento de Hanbury-Brown y 
Twiss (HBT) mide la probabilidad de detectar dos fotones con un retardo 
$\tau$. Veamos primero la evoluci\'on de la funci\'on de onda: 
%
\begin{eqnarray*}  	%\label{ec:G2operacional}
|\Psi(t+\tau) \rangle &=& A_{ge} (t+\tau) A_{ge} (t) |\Psi (0) \rangle  	\\
	&=& e^{iH(t+\tau)} A_{ge} e^{-iH(t+\tau)} e^{iHt} A_{ge} e^{-iHt} 
		|\Psi (0) \rangle \\ 
&=& e^{iH(t+\tau)} A_{ge} e^{-iH(t+\tau)} e^{iHt} A_{ge} |\Psi (t) \rangle \\ 
&=& e^{iH(t+\tau)} A_{ge} e^{-iH(t+\tau)} e^{iHt} |g \rangle \\ 
&=& e^{iH(t+\tau)} A_{ge} e^{-iH\tau}  |\Psi(t) \rangle \\ 
&=& e^{iH(t+\tau)} A_{ge} |\Psi(t+\tau) \rangle 
	= e^{iH(t+\tau)} |g \rangle \,.
\end{eqnarray*} 	
%%
Entonces, la correlaci\'on de emisi\'on de dos fotones es 
%
\begin{eqnarray*}  	\label{ec:G2operacional}
G^{(2)} (t,t+\tau) &=& \left| A_{ge} (t+\tau) A_{ge} (t) |\Psi (0) \rangle  \right|^2 	
	\nonumber \\ 
&=& \langle \Psi(0) | A_{eg} (t) A_{eg} (t+\tau)  \\ 
	&& A_{ge} (t+\tau) A_{ge} (t) |\Psi(0) \rangle \\
&=& \langle  A_{eg} (t) A_{eg} (t+\tau)  \\ 
	&& A_{ge} (t+\tau) A_{ge} (t) |\Psi(0) \rangle \langle \Psi(0) |  \rangle \\
&=& \langle  A_{eg} (t) A_{eg} (t+\tau)  A_{ge} (t+\tau) A_{ge} (t)\rho(0) \rangle \\
&=& \langle  A_{eg} (t) A_{eg} (t+\tau)  A_{ge} (t+\tau) A_{ge} (t) \rangle \,.
\end{eqnarray*} 	
%%

La detecci\'on de un fot\'on viene asociada al decaimiento del \'atomo por 
emisi\'on espont\'anea, por lo que el \'atomo pasa del estado excitado al 
estado base, descrito por el operador $A_{0j}$, perdiendo la memoria de la 
historia previa; estad\'isticamente, sin embargo, todas las historias son 
equivalentes si las condiciones iniciales se repiten. De este modo el sistema 
alcanza el estado estacionario, en el cual las funciones a dos tiempos se 
reducen a funciones del intervalo entre ellos, 
$G^{(2)} (t,t+\tau) \to G^{(2)} (\tau)$,  donde 
%
\begin{eqnarray}
G^{(2)} (\tau) 
	= \langle A_{eg} (0) A_{eg} (\tau)  A_{ge} (\tau) A_{ge} (0) \rangle \,.
\end{eqnarray} 	
%%
Cuando $\tau \to \infty$ los fotones emitidos son completamente 
independientes, por lo que 
%
\begin{eqnarray}
\langle A_{eg} (0) A_{eg} (\tau)  A_{ge} (\tau) A_{ge} (0) \rangle 
	&\to&  \langle A_{ee} \rangle(0) \langle A_{ee} \rangle(\tau) \nonumber \\ 
	&=& \langle A_{ee} \rangle_{st}^2 \,, 
\end{eqnarray} 	
%%
lo cual nos sirve para normalizar la correlaci\'on,
%
\begin{eqnarray}
g^{(2)} (\tau) 
	&=& \frac{ \langle A_{eg} (0) A_{eg} (\tau)  A_{ge} (\tau) A_{ge} (0) \rangle }
	{ \langle A_{ee} \rangle_{st}^2 } \,, 
\end{eqnarray} 	
%%
tal que  
%
\begin{eqnarray} 
g^{(2)} (\tau \to \infty) = 1 \,. 
\end{eqnarray} 	 
%%
%
\begin{eqnarray}
G^{(2)} (\tau) 
	&=& \langle A_{eg} (0) A_{eg} (\tau)  A_{ge} (\tau) A_{ge} (0) \rangle 	
		\nonumber \\
	&=& \langle A_{ee} \rangle_{st} 
	 \langle A_{eg} (\tau)  A_{ge} (\tau) \rangle_{\rho(0)=\rho_{gg}} 
	 	\nonumber\\ 
	&=& \langle A_{ee} \rangle_{st}  
	\langle A_{ee} (\tau) \rangle_{\rho(0)=\rho_{gg}} \,.
\end{eqnarray} 	
%%

Usando la f\'ormula cu\'antica de regresi\'on la correlaci\'on $g_j^{(2)}$ se 
puede reducir a la evoluci\'on de la poblaci\'on del estado excitado dado que 
inici\'o en el estado base, normalizada por la poblaci\'on promedio 
\cite{Carmichael02}, 
%
\begin{eqnarray}
g^{(2)} (\tau) &=& \frac{ \langle A_{ee} (\tau)\rangle_{\rho(0)=\rho_{gg}} }
	{ \langle A_{ee} \rangle_{st} }
\end{eqnarray}
%%
es decir, la correlaci\'on se reduce a normalizar la poblaci\'on del estado 
excitado. 

\textcolor{red}{Antibunching}

%%%%%%%%%%%%%%%%%%%%%%%%%%%%%%%%%%%
\subsection{Distribuci\'on de Tiempos de Espera o Funci\'on de Retardo} 
Una funci\'on de gran importancia en la fotoestad\'istica, y en muchos 
campos de la ciencia, es la distribuci\'on de intervalos entre dos fotones 
o eventos consecutivos. Mostramos los casos de un campo en estado 
coherente y de la fluorescencia resonante de un \'atomo de dos niveles. 

\subsubsection{Campo en Estado Coherente, M\'etodo 1 
\cite{Carmichael93} }
Procediendo como en la Secci\'on \ref{sec:DPF1} dividimos el tiempo 
de espera en $N$ subintervalos de duraci\'on $\Delta t$. En el l\'imite 
$N \to \infty, \, \Delta t \to 0$, con $N \Delta t =\tau$ constante, 
$w(\tau) = p q^N$, es decir, tras $N$ intervalos $q$ sin la emisi\'on 
de un fot\'on y luego la emisi\'on de uno, tenemos 
%
\begin{eqnarray}
w(\tau) d\tau = ( 1-\xi \bar{I} \Delta t )^N \xi \bar{I} \Delta t 
	\to \xi \bar{I}  e^{ -\xi \bar{I} \tau} d\tau \,,
\end{eqnarray}
%%
donde $\xi \bar{I} e^{ -\xi \bar{I} \tau}$ es la distribuci\'on de Poisson 
para $n=1$ y $\bar{n} = \xi \bar{I}$. El tiempo de espera promedio 
viene dado por 
%
\begin{eqnarray}
\bar{\tau} = \int_0^{\infty} d\tau \,\tau \, \xi \bar{I} e^{ -\xi  \bar{I} \tau} 
	= \frac{1}{\xi \bar{I} } \,.
\end{eqnarray}
%%
Este es justo el cociente del tiempo de conteo $T$ y el n\'umero 
promedio de fotoelectrones contados 
$\bar{n}$, $T/\bar{n} = T/ \xi \bar{I} T = 1/ \xi \bar{I}$, que es lo que 
esperar\'iamos. 

%%%%%%%%
\subsubsection{Campo en Estado Coherente, M\'etodo 2 \cite{KiKn89}} 
El campo coherente posee una distribuci\'on de Poisson de los fotones,  
$P(n) = \bar{n} e^{-\bar{n}} /n!$, donde $\bar{n}$ es n\'umero promedio 
de fotones. Si $\tau_0$ es el intervalo medio entre fotoconteos, habr\'a 
un n\'umero promedio $\tau/\tau_0$ de fotones en el intervalo $\tau$. 
La probabilidad $P(n, \tau)$ de $n$ fotones presentes en el campo en 
alg\'un intervalo $\tau$ es 
%
\begin{eqnarray}
P(n,\tau) = \frac{ (\tau/\tau_0)^n }{n!} e^{ -\tau/\tau_0} \,.
\end{eqnarray}
%%
Entonces la probabilidad $P(\tau)$ de no detectar ning\'un fot\'on en el 
retardo $\tau$ es simplemente
%
\begin{eqnarray}
P(\tau) \equiv P(0,\tau) = e^{ -\tau/\tau_0} \,.
\end{eqnarray}
%%
La integral de $w(\tau')$ con respecto a $\tau'$  sobre $[0,\tau]$ da la 
probabilidad de que el siguiente fot\'on sea emitido entre el tiempo de 
retraso 0 y $\tau$. El complemento de esta integral es la probabilidad de que 
ninguna emisi\'on haya occurrido de modo que 
%
\begin{eqnarray}
P(\tau) = 1 - \int_0^{\tau} w (\tau') \,d\tau \,,
\end{eqnarray}
%%
es decir, 
%
\begin{eqnarray}
w(\tau) = -dP/d\tau \,.
\end{eqnarray}
%%
Usando esta ecuaci\'on, la funci\'on de retardo (Loudon) es 
%
\begin{eqnarray}
w(\tau) = \frac{1}{\tau_0} e^{-\tau/\tau_0} \,.
\end{eqnarray}
%%
Esto se compara con $P(1,\tau) = (\tau/\tau_0) \exp{(-\tau/\tau_0)}$, 
que es la probabilidad de detectar un fot\'on en el intervalo de tiempo 
$\tau$. \textcolor{blue}{Pendiente: Ver \cite{CSVR89} p. 1204 para una 
derivaci\'on simple de $w_{coh}(\tau)$ usando $w_{2LA}(\tau)$. }

\textcolor{blue}{Pendiente: Este resultado es relevante en el estudio de 
fluorescencia intermitente, a interpretarse como la distribuci\'on de 
duraciones de per\'iodos oscuros. }

\subsubsection{Fluorescencia Resonante de un \'Atomo de dos Niveles 
	\label{sssec:wtd2LA}} %%%%%%%%%%%%%%%
\tbn{Carmichael, SMQO1 2.6} \cite{Carmichael02} (a) Demostrar que la 
distribuci\'on de tiempos de espera de fotoelectrones de fluorescencia 
resonante (\'atomo de dos niveles) para eficiencia de detecci\'on unidad 
esta dada por 
%
\begin{subequations}
\begin{eqnarray} 	
w_{ss}(\tau) &=& \frac{\gamma Y^2}{2Y^2-1} e^{-\gamma \tau/2}
	(1- \cosh{\delta' \tau}) \,, 	\label{eq:wtd_2LA} \\
%
\delta' &\equiv& (\gamma/2) \sqrt{1-2Y^2} \,, 	\quad  
	Y = \sqrt{2} \Omega/\gamma \,. \label{eq:delta_and_Y} \qquad 
\end{eqnarray}
\end{subequations}
%%
(b) Verificar que se satisface $\int_0^{\infty} w_{ss}(\tau) d\tau =1$. 
(c) Demostrar que el intervalo medio entre fotopulsos es 
%
\begin{eqnarray}
\overline{\tau} =2(1+Y^2)/\gamma Y^2 
=(\gamma \langle \sigma_+ \sigma_- \rangle_{ss} )^{-1} \,.
\end{eqnarray}
%%
(d) Graficar $w_{ss}(\tau)$ para $2Y=1$ y comparar con la exponencial 
$w_{ss}(\tau) = (\gamma/6) e^{-\gamma \tau /6}$ que se obtiene para luz 
coherente de la misma intensidad. \textit{Nota:} Esta f\'ormula para 
$w_{ss}(\tau)$ esta limitada al caso de frecuencias iguales de l\'aser y de 
la transici\'on at\'omica. Para el caso no resonante utilizar la soluci\'on 
general para $c_e(\tau)$ de la ecuaci\'on de Schr\"odinger con el 
Hamiltoniano efectivo, sin colapso por emisi\'on de un fot\'on. 
\textcolor{blue}{Tarea Extra: Obtener 
$\Delta \tau^2 = \overline{\tau^2} -\overline{\tau}^2$ usando Mathematica.}
%
\begin{figure}[b]
\includegraphics[width=8.5cm,height=6.25cm]{fig1_wtd.pdf} 
\caption{\label{wtd} 
Funci\'on de retardo para la fluorescencia de un \'atomo de dos niveles 
(azul) y luz coherente (naranja), ambas de la misma intensidad. Inciso 
(d) del Ejercicio 2.6 en \cite{Carmichael02}:} 
\end{figure}
%% 

\textit{Demostraciones}
\newline (a) La funci\'on buscada est\'a dada por 
$w(\tau) = \gamma |c_e(\tau)|^2$, donde $c_e(\tau)$ se obtiene de 
resolver las ecuaciones de la trayectoria sin ocurrir una emisi\'on, 
Ecs.(\ref{eq:QJ2LA-sol}) abajo y (8.22b) de \cite{Carmichael93}:
%
\begin{eqnarray*}
w_{ss}(\tau) &=& \gamma \left|  \frac{\Omega}{2\delta} e^{-\gamma \tau/4} 
	\sinh{\delta \tau} \right|^2 \,, 	\quad 
	2\delta = \frac{1}{2} \sqrt{\gamma^2 -4\Omega^2} \,, \\
%
&=& \frac{\gamma 2\Omega^2}{(\gamma^2 -4\Omega^2)/2} 
	e^{-\gamma \tau/2} \frac{1}{2} \left( \cosh{2\delta \tau}  -1 \right) \\
&=&  \frac{\gamma 2\Omega^2}{\gamma^2 (\frac{4\Omega^2}{\gamma^2} -1)}  
	e^{-\gamma \tau/2}  \left( 1- \cosh{2\delta \tau}  \right)  \,. 
\end{eqnarray*}
%%
Con $2\delta = \delta'$ en (\ref{eq:delta_and_Y}) obtenemos 
(\ref{eq:wtd_2LA}).  

La funci\'on de retardo muestra la propiedad de \textit{antibunching}, 
que expresa que el \'atomo no emite dos fotones a la vez, es decir, 
$w(\tau) =0$. 

(b) Normalizaci\'on de la funci\'on de retardo: 
\begin{widetext}
%
\begin{eqnarray*}
\int_0^{\infty} w_{ss}(\tau) d\tau 
&=& \frac{\gamma Y^2}{2Y^2-1} \int_0^{\infty} d\tau \, e^{-\gamma \tau/2} 
\left[ 1- \frac{ e^{\delta' \tau} +e^{-\delta' \tau} }{2} \right] 
=\frac{\gamma Y^2}{2Y^2-1} \left[ \frac{e^{-\gamma \tau/2}}{-\gamma/2} 
	-\frac{ e^{(\delta' -\gamma/2)\tau}}{2(\delta' -\gamma/2)} 
	-\frac{ e^{-(\delta' +\gamma/2)\tau}}{-2(\delta' +\gamma/2)}
	\right]_0^{\infty} 	\\
&=& \frac{\gamma Y^2}{2Y^2-1} \left[ \frac{2}{\gamma} 
	+\frac{1}{2(\delta' -\gamma/2)} -\frac{1}{-2(\delta' +\gamma/2)}
	\right]		
=\frac{\gamma Y^2}{2Y^2-1} \left[ \frac{2}{\gamma} 
	+\frac{\delta' +\gamma/2 -\delta' +\gamma/2}
	{2(\delta'^2 -(\gamma/2)^2)}  \right]		\\
&=& \frac{\gamma Y^2}{2Y^2-1} \left[ \frac{2}{\gamma} 
	+\frac{\gamma/2}{(\gamma/2)^2 -\Omega^2 -(\gamma/2)^2)} \right] 
=\frac{\gamma Y^2}{2Y^2-1} \left[ \frac{2}{\gamma} 
	-\frac{\gamma}{2\Omega^2} \right] 
=\frac{\gamma Y^2}{2Y^2-1} \left[ 
	\frac{4\Omega^2 -\gamma^2}{2\Omega^2 \gamma} \right] 	\\
&=& \frac{2\gamma \Omega^2/\gamma^2}{4(\Omega^2/\gamma^2) -1} 
  \left[ \frac{4\Omega^2 -\gamma^2}{2\Omega^2 \gamma} \right] 	
=1 \qquad \qed 
\end{eqnarray*}
%%  
(c) Intervalo medio entre fotopulsos:
%
\begin{eqnarray*}
\overline{\tau} &=& \int_0^{\infty} \tau w_{ss}(\tau) d\tau 
= \frac{\gamma Y^2}{2Y^2-1} \int_0^{\infty} d\tau \,\tau \, e^{-\gamma \tau/2} 
\left[ 1- \frac{ e^{\delta' \tau} +e^{-\delta' \tau} }{2} \right] \,,
%
\qquad \qquad \int d\tau\,\tau \,e^{a\tau} 
=\frac{\tau \,e^{a\tau}}{a} -\frac{e^{a\tau}}{a^2} \,, 	\\
&=& \frac{\gamma Y^2}{2Y^2-1} \left[ \frac{\tau\,e^{-\gamma \tau/2}}{-\gamma/2} 
	-\frac{e^{-\gamma \tau/2}}{(-\gamma/2)^2} 
	-\frac{ \tau\,e^{(\delta' -\gamma/2)\tau}}{2(\delta' -\gamma/2)} 
	+\frac{ e^{(\delta' -\gamma/2)\tau}}{2(\delta' -\gamma/2)^2} 
	-\frac{ \tau\,e^{-(\delta' +\gamma/2)\tau}}{-2(\delta' +\gamma/2)}
	+\frac{ e^{-(\delta' +\gamma/2)\tau}}{2(\delta' +\gamma/2)^2}
	\right]_0^{\infty} 	\\
&=& \frac{\gamma Y^2}{2Y^2-1} \left[ \frac{1}{(\gamma/2)^2} 
	-\frac{1}{2(\delta' -\gamma/2)^2} 
	-\frac{1}{2(\delta' +\gamma/2)^2} \right] 
=\frac{\gamma Y^2}{2Y^2-1} \left\{ \frac{1}{(\gamma/2)^2} -\frac{1}{2} 
\left[ \frac{(\delta' +\gamma/2)^2 +(\delta' -\gamma/2)^2}
	{(\delta'^2 -(\gamma/2)^2)^2} 	\right] \right\} 	\\
&=& \frac{\gamma Y^2}{2Y^2-1} \left[ \frac{1}{(\gamma/2)^2} 
   -\frac{\delta'^2 +(\gamma/2)^2}{(\delta'^2 -(\gamma/2)^2)^2} \right] 
=\frac{\gamma Y^2}{2Y^2-1} \left[ \frac{1}{(\gamma/2)^2} 
   -\frac{(\gamma/2)^2 (1-2Y^2) +(\gamma/2)^2}
	{[(\gamma/2)^2 (1-2Y^2) -(\gamma/2)^2]^2} \right] 	\\
&=& \frac{\gamma Y^2(\gamma/2)^2}{2Y^2-1} 
	\left[ \frac{1}{(\gamma/2)^4} -\frac{2(1-Y^2)}{4Y^4} \right] 
=\frac{\gamma Y^2(\gamma/2)^2}{2Y^2-1} 
	\left[ \frac{2Y^4 -1 +Y^2}{2Y^4(\gamma/2)^4} \right] 
=\frac{\gamma Y^2}{2Y^2-1} 
	\left[ \frac{(2Y^2 -1)(1 +Y^2)}{2Y^4(\gamma/2)^2} \right]  \\
\overline{\tau} &=& \frac{2(1 +Y^2)}{\gamma Y^2} \,, 
	\qquad \qed 
\end{eqnarray*}
%%    
\vspace{2cm}
\end{widetext}
% 
%\textcolor{red}{Obtener este resultado mediante 
%$\overline{\tau} = (\gamma \rho_{ee})^{-1}$.}
\textcolor{blue}{Tarea: Obtener $w(\tau)$ para el caso de shelving.}

%%%%%%%%10%%%%%%%%20%%%%%%%%30%%%%%%%%40%%%%%%%%50%%%%%%%%60%%%%%%%%70%%%
\section{Trayectorias Cu\'anticas: Introducci\'on}
These notes are meant as a compilation of the main QT approaches to 
RF for use in class. The notes are limited to the very basics (of the physics 
and theory of QTs) but hopefully the students will be able to develop their 
own computer codes. 

Intensity, intensity-intensity correlations, waiting-time distribution, 
photoelectron statistics, and sub-Poissonian statistics are measured by 
direct photon counting, simulated by discontinuous trajectories (QJ). 
Homodyne and heterodyne detection, simulated by quantum state 
diffusion trajectories, allow phase-dependent field measurements, such 
as single-quadrature, squeezing, and spectra. Conditional Homodyne 
Detection (CHD) mixes QJs and QSD to provide a new look at quantum 
fluctuations and resolves squeezing in a detector-efficiency-independent 
manner. Spectra can be measured by homodyne, heterodyne, and filtered 
measurements \cite{TiCa92,Holland98}. 

Quantum trajectories, as formulated by Carmichael, simulate the output 
of an experiment, and a large ensemble of trajectories represents the 
average result of many experiments. 


%%%%%%%%%%%%%%%%%%%%%%%%%%%%%%%
\subsection{Evolution of an Open Quantum System} 
Our system consists of a coherent part that features the atom-laser 
coupling. This is often studied with a unitary evolution via the 
Schr\"odinger equation. However, the coupling of the atom to the 
environment leads to dissipation, thus forcing a density operator 
approach (master equation) or a Heisenberg-Langevin approach. 

In resonance fluorescence, we have an atomic system driven by lasers 
and maybe cavity fields. Magnetic fields are often added to prepare 
the atomic system in a specific state. To simplify matters, we mostly speak 
of the interaction of a single two-level atom of transition frequency 
$\omega_a$ with a single monochromatic laser treated classically (not an 
operator). The atom-laser coupling is given by the Rabi frequency 
$\Omega$. The (RWA) Hamiltonian in the frame rotating at the laser 
frequency is 
%
\begin{equation}
H_0 = \Delta \sigma_{ee} +\frac{\Omega}{2}(\sigma_+  +\sigma_-)  \,,
\end{equation}
%%
where $\sigma_i$ are Pauli spin operators, and 
$\Delta= \omega_a -\omega$ is the atom-field detuning ($\hbar=1$). 
The state of the atom is given by the wave function 
$\ket{\Psi (t)} = c_g(t) \ket{g} +c_e(t) \ket{e}$. 

The atom is also coupled to a large reservoir of harmonic oscillators,  
which induce decay of the excited state at the rate $\gamma$. We can 
neglect blackbody radiation effects at optical frequencies for room 
temperatures and below. The master equation describes the full 
atom-laser-reservoir 
%
\begin{eqnarray} 	\label{MasterEq}
\dot{\rho} &=& - i [H_0,\rho] +\frac{\gamma}{2} \left[ 2 \sigma_- \rho \sigma_+ 
	-\sigma_+ \sigma_- \rho - \rho  \sigma_+ \sigma_- \right]  	\,. 	\qquad	
\end{eqnarray}
%%

The density matrix formalism ignores the shot noise or, rather, averages 
over an infinite number of realizations.

%%%%%%%%10%%%%%%%%20%%%%%%%%30%%%%%%%%40
%%%%%%%%50%%%%%%%%60%%%%%%%%70%%%
\section{Trayectorias Discont\'inuas: Simulaci\'on de Fotoconteo}
En este caso la luz se recolecta por detecci\'on directa de los fotones por 
fotodetectores de banda ancha. Escribimos la ecuaci\'on maestra como 
%
\begin{eqnarray} 	\label{MasterEqEffJ}
\dot{\rho} &=& \mcL \rho = - i ( \mcH \rho -\rho \mcH^{\dag} )  +\mcS \rho \,, 
\end{eqnarray}
%%
donde a $\mcL$ se le llama superoperador de Lindblad, 
%
\begin{eqnarray} 
\mcH &=& H_0 - i \frac{\gamma}{2} \sigma_+ \sigma_-   
\end{eqnarray}
%%
es un Hamiltoniano efectivo no Hermitiano, y  
%
\begin{eqnarray} 
\mcS \rho &=& \gamma \sigma_- \rho \sigma_+ 
\end{eqnarray}
%%
es llamado \textit{el operador de salto}, que corresponde a un evento de 
emisi\'on espont\'anea. Esto es porque  
$\gamma \sigma _- \rho \sigma_ + = \gamma \rho_{ee} \sigma_{gg}$, es 
decir,  el operador de salto disminuye el estado del \'atomo (reduce la 
funci\'on de onda). 

The master equation can be unraveled in infinite ways, but only a limited 
number have an experimental significance, each looking at different 
properties. The decomposition above is suitable for direct photodetection 
simulation, it corresponds to a discontinuous evolution of the wave 
function due to photon emissions (quantum jumps); this deals with the 
particle aspect of light  \cite{Carmichael93,Carmichael02,DaCM92,DuZR92}. We later deal 
with homodyne and heterodyne measurements (wave features -- 
quantum diffusion) 
\cite{Carmichael93,Carmichael02,WiMi93a,WiMi93b,Percival98}. 
Also spectral detection \cite{TiCa92,Holland98}. 

La soluci\'on formal  de la ecuaci\'on maestra es $\rho(t) = e^{\mcL t} \rho (0)$. 
Con la descomposici\'on (\ref{MasterEqEffJ}) podemos escribir 
%
\begin{eqnarray} 	%\label{MasterEqEffJ}
\dot{\rho} &=& ( \mcL - \mcS ) \rho  +\mcS \rho \,. 		
\end{eqnarray}
%%
La soluci\'on 
%
\begin{eqnarray} 	%\label{MasterEqEffJ}
\rho (t) &=& e^{ [( \mcL - \mcS ) +\mcS ] t} \rho (0)  	\nonumber \\ 
&=& 	\sum_{m=0}^{\infty} \int_0^t dt_m  \int_0^{t_m} dt_{m-1}  \cdots 	
	 \int_0^{t_2} dt_1 e^{ (\mcL - \mcS) (t-t_m)}  		\nonumber \\ 
	&& \mcS e^{ (\mcL - \mcS) (t_m-t_{m-1})} \mcS \cdots 
	\mcS  e^{ (\mcL - \mcS) t_1} \rho(0) \,,
\end{eqnarray}
%%
no es dif\'icil de entender \cite{Carmichael93}. El integrando es el 
operator de densidad condicional no normalizado  $\bar{\rho}_c (t)$ 
%
\begin{eqnarray} 	\label{eq:condDensOp}
\bar{\rho}_c (t) &=& e^{ (\mcL - \mcS) (t-t_m)} \mcS 	
	 e^{ (\mcL - \mcS) (t_m-t_{m-1}} \mcS 	\nonumber \\  
	&&\cdots \mcS  e^{ (\mcL - \mcS) t_1} \rho(0) \,,
\end{eqnarray}
%%
para el estado inicial $\rho_c (0) = \rho (0)$. Esta serie es una suma sobre diferentes 
historias de fotoemisiones de la fuente desde $t=0$ hasta $t$. Cada historia es 
\'unica y puede involucrar cualquier n\'umero de fotoemisiones, desde $m=0$ hasta $m=\infty$, y los tiempos de las emisiones pueden ser cualquier secuencia ordenada 
de $m$ veces en el intervalo $[0,t]$.  La interpretaci\'on f\'isica de $\bar{\rho}_c (t)$ 
es de evoluciones sin fotoemisiones interrumpidas por \textit{colapsos} del estado a 
los tiempos de emisi\'on. Al tiempo $t$, para un estado inicial $\rho (0)$ y una 
secuencia particular de tiempos de emisi\'on, el operador de densidad condicionado 
de la fuente esta dado por 
%
\begin{eqnarray} 	\label{eq:condDensOp}
\rho_c (t) &=& \frac{\bar{\rho}_c (t) }{ \Tr [\bar{\rho}_c (t) ]} \,.
\end{eqnarray}
%%

Cada secuencia o historia de fotoemisiones se llama \textit{trayectoria cu\'antica}. 
Este tipo de trayectorias de fotoemisiones (detecci\'on directa)  se denomina en 
ingl\'es \textit{unraveling}, una descomposici\'on particular del operador de densidad. 
Existen otras descomposiciones que corresponden a otros m\'etodos de medici\'on 
de la luz. 

N\'otese que mientras que $\rho (t)$ llega a un  estado estacionario, $\rho_c (t)$ 
no lo har\'a; cada trayectoria mantiene su evoluci\'on estoc\'astica. La trayectoria, 
sin embargo, es estacionaria en sentido estad\'istico, y el estado estacionario 
para $\rho$ puede calcularse como un promedio en el tiempo de $\rho_c (t)$. 

Para muchos sistemas el Hamiltoniano permite factorizar el operador 
de densidad en \textcolor{red}{estados puros}
%
\begin{eqnarray}
\rho_c (t) &=& \ketbra{\Psi_c (t) }{ \Psi_c  (t)} 	\,, 	\quad 
	\bar{\rho}_c (t) = \ketbra{ \bar{\Psi}_c (t) }{ \bar{\Psi}_c  (t)} \,, 
	\nonumber \\ 
\end{eqnarray}
%%
Mas adelante veremos algunos ejemplos. Entonces reemplazamos el propagador 
$e^{ (\mcL - \mcS) \Delta t}$ del operador de densidad $\bar{\rho}_c (t)$ por el de 
la funci\'on de onda $| \bar{\Psi}_c (t) \rangle$. La propagaci\'on sin fotoemisi\'on 
en un tiempo $\Delta t$ esta dada por 
%
\begin{eqnarray}
\ket{ \bar{\Psi}_c (t+ \Delta t) } 
	= e^{- i \mcH \Delta t /\hbar} \ket{ \bar{\Psi}_c (t) } 
\end{eqnarray}
%%
que es la soluci\'on de la ecuaci\'on de Schr\"odinger 
%
\begin{eqnarray}
i \hbar \frac{d}{dt} \ket{ \bar{\Psi}_c } = {\cal H} \ket{ \bar{\Psi}_c } \,,
\end{eqnarray}
%%
donde $\ket{ \bar{\Psi}_c }$ es una funci\'on de onda no-normalizada 
debida a la acci\'on del Hamiltoniano no Hermitiano $\mcH$ que lleva 
a una evoluci\'on no unitaria, asegurando la p\'erdida de probabilidad 
de que el \'atomo se encuentre en el estado excitado. 

Escribamos la funci\'on de onda no normalizada como 
%
\begin{equation} 	\label{eq:2LA-WF-unnorm}
\ket{ \bar{\Psi}_c (t) } = \bar{c}_g(t) \ket{g} +\bar{c}_e(t) \ket{e} 	\,,
\end{equation}
%%
La funci\'on de onda normalizada es 
%
\begin{equation} 	
\ket{ \Psi_c(t) } = \ket{ \bar{\Psi}_c(t) } / 
	[\braket{ \bar{\Psi}_c(t)}{ \bar{\Psi}_c(t)} ]^{1/2} \,.
\end{equation}
%%
El sistema \'atomo-campo evoluciona de manera cont\'inua pero no 
unitaria por un per\'iodo de tiempo hasta que se emite un fot\'on 
espont\'aneamente, se\~nalando el regreso del \'atomo al estado 
base. Este ciclo se repite durante una historia. 

Reescribamos el superoperador de brinco $\mcS \rho$ como 
%
\begin{eqnarray} 
\mcS \rho &=&  \mcC \rho \mcC^\dag \,, 	\qquad \to \qquad 
	\mcC = \sqrt{\gamma} \sigma_- 	\,. 
\end{eqnarray}
%%
Entonces, al tiempo $t_j$ en el que ocurre una emisi\'on el estado no 
normalizado experimenta un colapso dado por el operador de brinco 
$\mcC$,  
%
\begin{equation}
\ket{ \bar{\Psi}_c } \to \mcC \ket{ \bar{\Psi}_c } 
	= \sqrt{\gamma}\, \bar{c}_e (t_j) \ket{g} 	\,,
\end{equation}
%%
que se normaliza como 
 %
\begin{eqnarray} 	\label{jump}
\ket{\Psi_c(t_j)} = \frac{\mcC \ket{\Psi_c(t_j) } }
  	 {\sqrt{\bra{\Psi_c(t_j)}  \mcC^\dag \mcC \ket{\Psi_c(t_j)} } } \,,	
\end{eqnarray}
%%
y que ocurre con probabilidad en el intervalo $[t, t+ \delta t]$ 
%
\begin{eqnarray} 	\label{JumpProb}
p_c(t) &=& \delta t \bra{\Psi_c (t)} \mcC^\dag \mcC \ket{\Psi_c (t)} 	\\ 
	&=& \gamma \delta t \bra{\Psi_c (t)} \sigma_+ \sigma_- \ket{\Psi_c (t)}
	\nonumber 	\\ 
&=& \gamma\, \delta t \, \frac{|\bar{c}_e(t)|^2}
	{|\bar{c}_g(t)|^2 +|\bar{c}_e(t)|^2} = \gamma\, \delta t |c_e(t)|^2	\,. 	
	\qquad 	\label{eq:jumpProb2LA}
\end{eqnarray}
%%

Entonces, una simulaci\'on num\'erica ocurre en el tiempo discreto con 
paso $\delta t$. Obtenemos una trayectoria estoc\'astica de la funci\'on 
de onda $\ket{\Psi_c (t_n)}$, donde $t_n = n \delta t$. Dada 
$\ket{\Psi_c (t_n)}$, la funci\'on de onda $\ket{\Psi_c (t_{n+1})}$ se obtiene 
con el siguiente algoritmo: 
\newline (i) Evaluar la probabilidad de un brinco $p_c (t_n)$. 
\newline (ii) Generar un n\'umero aleatorio $r_n$ distribuido uniformemente en el intervalo $[0,1]$. 
\newline (iii) Comparar $p_c(t_n)$ con $r_n$ y calcular  $\ket{\Psi_c (t_{n+1})}$ 
de acuerdo a la regla   
 %
\begin{subequations}
\begin{eqnarray} 	%\label{jump}
\ket{\Psi_c(t_{n+1}) } &=& \frac{\mcC \ket{\Psi_c(t_n) } } 
	{\sqrt{\bra{\Psi_c(t_n)} \mcC^\dag \mcC \ket{\Psi_c(t_n)} } } \,, 		
	\quad p_c (t_n) > r_n \,, 	\nonumber \\ \\ 
\ket{\Psi_c(t_{n+1}) } 
	&=& \frac{ e^{- i \mcH \delta t /\hbar} \ket{ \Psi_c(t_n) } } 
	{\sqrt{\bra{\Psi_c(t_n)} e^{ i \mcH \delta t /\hbar} 
	e^{- i \mcH \delta t /\hbar} \ket{\Psi_c(t_n)} } } 	\,,	\nonumber \\ 
	&& \qquad  \qquad  \qquad  \qquad p_c (t_n) \leq r_n \,. 	
\end{eqnarray}
\end{subequations}
%%
El resultado es un mapeo estoc\'astico cu\'antico a los tiempos $t_n$ donde  
los brincos ocurren separados por muchos intervalos $\delta t$: 
 %
\begin{eqnarray} 	
\lefteqn{ \ket{\Psi_c(t_{n+1}) } =} \hspace{1cm} \nonumber \\
	&& \frac{ \mcC e^{- i \mcH \tau_{n+1} /\hbar} \ket{\Psi_c(t_n) } } 
	{\sqrt{\bra{\Psi_c(t_n)} e^{ i \mcH \tau_{n+1} /\hbar} \mcC^\dag 		
	\mcC e^{- i \mcH \tau_{n+1} /\hbar} \ket{\Psi_c(t_n)} } } 	\,,	
	\nonumber \\
\end{eqnarray}
%%
donde $\tau_{n+1} = t_{n+1} -t_n$ es un tiempo aleatorio cuya 
estad\'istica depende de la funci\'on estoc\'astica misma mediante 
$p_c (t_n)$. El algoritmo incorpora esta estad\'istica ``al vuelo'' al 
tomar muchos pasos infinitesimales $\delta t$. 

%%%%%%%%%%%%%%%%%%%%%%%%%%%%%%%
\subsection{Ejemplos}
\subsubsection{Emisi\'on Espont\'anea} %%%%%%%%%%
%\textcolor{blue}{Tarea}

%\subsubsection{Decaimiento de un Campo de una Cavidad} %%%%%
\subsubsection{Fluorescencia Resonante} %%%%%%%%%%%%%
Para un \'atomo de dos niveles (A2N) con funci\'on de onda 
Ec.(\ref{eq:2LA-WF-unnorm}) interactuando con un l\'aser 
monocrom\'atico en espacio libre el Hamiltoniano efectivo es 
%
\begin{eqnarray} 
\mcH &=& (\Delta - i \frac{\gamma}{2})  \sigma_+ \sigma_-   
	+\frac{\Omega}{2} (\sigma_+ +\sigma_-) \,,
\end{eqnarray}
%%
y el operador de colapso es $ \mcC = \sqrt{\gamma} \sigma_-$. Las 
ecuaciones para las amplitudes son:
%
\begin{subequations}
\label{eq:DirectSSEamp}
\begin{eqnarray}
i \dot{\bar{c}}_g &=& \frac{\Omega}{2} \bar{c}_e  \,, \\
%
i \dot{\bar{c}}_e &=& \frac{\Omega}{2} \bar{c}_g 
	+(\Delta -i\frac{\gamma}{2}) \bar{c}_e \,,
\end{eqnarray}
\end{subequations}
%%
Podemos obtener soluciones anal\'iticas si $\Delta=0$ (v. gr., por 
transformada de Laplace). Con las condiciones iniciales $\bar{c}_g(0) =1$
y $\bar{c}_e(0)=0$, las soluciones son: 
%
\begin{subequations}		\label{eq:QJ2LA-sol}
\begin{eqnarray}
\bar{c}_g (t) &=& e^{-\gamma t/4} \left[ \cosh{ \delta t} 
	+\frac{\gamma/2}{2\delta} \sinh{\delta t}   \right]  \,, \\
%
\bar{c}_e (t) &=& i \frac{\Omega}{2\delta} e^{-\gamma t/4} \sinh{\delta t}  \,, \\
%
2\delta &=& \frac{1}{2} \sqrt{\gamma^2 -4\Omega^2}  \,.
\end{eqnarray}
\end{subequations}
%%
La probabilidad de un salto est\'a dada por la Ec.(\ref{eq:jumpProb2LA}). 

Una aplicaci\'on simple de implementar es la distribuci\'on de tiempos 
de espera, vista en \ref{sssec:wtd2LA}. 

%\textcolor{red}{El cero de energ\'ia est\'a a la mitad de la transici\'on, 
%por lo que las desinton\'ias estan divididas entre 2. Cambiar al caso de 
%energ\'ia cero para el estado base. }

%%%%%%%%%%%%%%%%%%%%%%%%%%%%%%%
\subsection{C\'alculo de Observables a un Tiempo}
El valor esperado de un operador $A$ al tiempo $t$ se obtiene como el 
promedio sobre un n\'umero muy grande $n$ de trayectorias 
 %
\begin{eqnarray} 	 
\vxp{A (t)} &=& 
	\frac{1}{n} \sum_{i=1}^n \bra{\Psi_c^i (t)} A \ket{\Psi_c^i (t)} \,.
\end{eqnarray}
%%

%%%%%%%%%%%%%%%%%%%%%%%%%%%%%%%
\subsection{Recuperando la Ecuaci\'on Maestra}
Aqu\'i demostramos que un ensamble infinito de trayectorias en el 
esquema de detecci\'on directa resuelve (unravels) la ecuaci\'on maestra. 
El cambio en $\ketbra{\Psi }{ \Psi}$ se debe a una parte de brinco m\'as 
una parte de evoluci\'on cont\'inua, ambas ponderadas apropiadamente. 
Para simplificar la tipograf\'ia usamos  $b^\dag =\sigma_+$ y 
$b= \sigma_-$, y $|\bar{\Psi}_c \rangle \to |\Psi \rangle$. Por integraci\'on 
formal de  $i(d/dt) \ket{\Psi(t)} = \mcH \ket{\Psi(t)}$ tenemos 
$\ket{\Psi(t+dt)} = (1 -i \mcH dt) \ket{\Psi(t)}$, donde $\ket{\Psi}$ es la 
funci\'on de onda condicionada no normalizada, 
%
\begin{equation}
d \rho = \ketbra{\Psi(t+dt) }{ \Psi(t+dt)} - \ketbra{\Psi(t) }{ \Psi(t)} 	\,,
\end{equation}
%%
\begin{widetext}
%
\begin{eqnarray*}
\ketbra{\Psi(t+dt) }{ \Psi(t+dt)}  
	& \rightarrow &	p_c \ketbra{\Psi (t) }{ \Psi (t)} 	
	 +(1 -p_c) \ketbra{\Psi_{no}(t) }{ \Psi_{no}(t)} 	\\
&=& 	\gamma dt \langle \Psi |b^{\dag}b| \Psi \rangle 
	\frac{b \ketbra{\Psi(t) }{ \Psi(t)} b^\dag}
	{ \bra{\Psi(t)} b^\dag b \ket{\Psi(t)} }		
	 +(1 -p_c)[1 - i \mcH dt]	
	\frac{ \ketbra{\Psi (t) }{ \Psi (t)} }{1-p_c}  [1 +i \mcH^\dag dt]	\\
&=& 	\gamma dt  \,b \ketbra{\Psi(t) }{ \Psi(t)} b^\dag 		
	 +[1 - i \mcH dt]	 \ketbra{\Psi (t) }{ \Psi (t)}   [1 +i \mcH^\dag dt]	\\
&=& \gamma dt\, b \ketbra{\Psi (t) }{ \Psi (t)} b^\dag +\ketbra{\Psi (t) }{ \Psi (t)}  		 
	+\frac{1}{i\hbar} dt [H_0, \ketbra{\Psi (t) }{ \Psi (t)} ]	\\ 	
	 && -\frac{\gamma}{2} dt \left[ b^\dag b \ketbra{\Psi (t) }{ \Psi (t)} 	 
	 +\ketbra{\Psi (t) }{ \Psi (t)} b^\dag b \right] 	+O(dt^2)  \,.
\end{eqnarray*}
%%
\end{widetext}
Restando $\ketbra{\Psi (t)}{ \Psi(t)}$ en ambos lados, dividiendo entre  $dt$  
e ignorando los t\'erminos $O(dt^2)$ se tiene 
%
\begin{eqnarray}
\frac{d \rho}{dt} &=& 
  \frac{ \ketbra{\Psi(t+dt) }{ \Psi(t+dt)} - \ketbra{\Psi (t) }{ \Psi (t)} }{dt}	
  \nonumber \\
	&=& \frac{1}{i\hbar} [H_0, \ketbra{\Psi (t) }{ \Psi (t)} ] 
	+\gamma \, b \ketbra{\Psi (t) }{ \Psi (t)} b^\dag	 	\nonumber\\ 
	&& -\frac{\gamma}{2} \left[ b^\dag b \ketbra{\Psi (t) }{ \Psi (t)} 
	+\ketbra{\Psi (t) }{ \Psi (t)} b^\dag b \right] 	  \,. \qquad 
\end{eqnarray}
%%
Promediando sobre un ensamble infinito de trayectorias, 
$ \overline{ \ketbra{\Psi (t) }{ \Psi (t)} }$, y haciendo 
$d \rho /dt \to \dot{\rho}$, reproducimos la ecuaci\'on maestra, 
Eq.~ (\ref{MasterEq}). 

Dado que la EES representa la evoluci\'on de un estado puro, una sola 
trayectoria conservar\'a la magnitud del vector de  Bloch, 
$\rho_c= |\Psi_c \rangle \langle \Psi_c|$, $\Tr (\rho_c)=1$.


%%%%%%%%10%%%%%%%%20%%%%%%%%30%%%%%%%%40
%%%%%%%%50%%%%%%%%60%%%%%%%%70%%%
\section{Detecci\'on Homodina Balanceada} 
(Seg\'un Scully \& Zubairy p.125. Revisar.)
\newline La detecci\'on de estados comprimidos requiere un esquema 
sensible a la fase que mide la varianza de una cuadratura del campo. 
Para ello se requiere interferir la luz del emisor con un campo EM de 
referencia que llamamos oscilador local (OL). En la detecci\'on homodina 
el OL tiene la misma frecuencia $\omega_{lo}$ e  intensidad es mucho 
mayor que el haz de excitaci\'on $\omega$. De hecho, ambos haces 
provienen del mismo l\'aser. Tambi\'en, el OL tiene un retraso de fase 
$\phi$, usualmente para empatar con una de las cuadraturas.  
%
\begin{figure}[t]
\includegraphics[width=9cm,height=6.cm]{fig_bhd.png} 
\caption{\label{dHB} 
Esquema de detecci\'on homodina y heterodina balanceada. En la 
DHm un mismo l\'aser impulsa al emisor (c\'irculo rojo) y al OL y se 
puede seleccionar la fase del OL (v. gr., $\phi =0,\pi/2$). En la DHt 
emisor y OL est\'an muy fuera de sinton\'ia y se pierde la selecci\'on 
de la fase. (fig. por Kevin MF)} 
\end{figure}
%% 

Sean $\hat{a}$ y $\hat{b}$ los modos del emisor y del OL, 
respectivamente. Ambos modos se mezclan mediante un divisor de haz 
(DH) con reflectancia $R$ y transmitancia $T$, tales que $R+T =1$.
Tras el DH emergen los modos 
%
\begin{eqnarray}
 \hat{c} = \sqrt{T} \hat{a} +i \sqrt{1-T} \hat{b} \,, 	\quad 
 \hat{d} = i \sqrt{1-T} \hat{a} + \sqrt{T} \hat{b} + \,, \,\,\,	
\end{eqnarray}
%% 
que llegan a los detectores  D1 y D2, respectivamente. Recu\'erdese 
que hay un cambio de fase $\pi/2$ ($e^{i \pi/2} =i$) entre las ondas 
transmitida y reflejada. Los modos $\hat{a}$ y $\hat{b}$ tienen la misma 
frecuencia para evitar introducir dependencia temporal a los operadores.  

Las intensidades registradas en D1 y D2 son 
%
\begin{eqnarray*}
\hat{c}^+ \hat{c} &=& T \hat{a}^+ \hat{a} +(1-T) \hat{b}^+ \hat{b} 
	+i \sqrt{T(1-T)}  \left( \hat{a}^+ \hat{b} -\hat{b}^+ \hat{a}  \right) \,, \\ 
\hat{d}^+ \hat{d} &=& (1-T) \hat{a}^+ \hat{a} +T \hat{b}^+ \hat{b} 
	-i \sqrt{T(1-T)}  \left( \hat{a}^+ \hat{b} -\hat{b}^+ \hat{a}  \right) \,. 
\end{eqnarray*}
%% 
Si $R=T=1/2$, 
%
\begin{subequations}
\begin{eqnarray}
\hat{c}^+ \hat{c} &=& \frac{1}{2} \hat{a}^+ \hat{a} + \hat{b}^+ \hat{b} 
	+\frac{i}{2} \left( \hat{a}^+ \hat{b} -\hat{b}^+ \hat{a}  \right) \,, \\ 
\hat{d}^+ \hat{d} &=& \frac{1}{2} \hat{a}^+ \hat{a} +\frac{1}{2} \hat{b}^+ \hat{b} 
	- \frac{i}{2} \left( \hat{a}^+ \hat{b} -\hat{b}^+ \hat{a}  \right) \,. 
\end{eqnarray}
\end{subequations}
%% 
En la detecci\'on homodina ordinaria solamente se mide uno de los 
modos ($\hat{c}$ \'o $\hat{d}$), mientras que en la Detecci\'on Homodina 
Balanceada (DHm) se restan las intensidades de D1 y D2, 
%
\begin{eqnarray}
\hat{n}_{cd} = \hat{c}^+ \hat{c} - \hat{d}^+ \hat{d} 
= - i \left( \hat{a}^+ \hat{b} -\hat{b}^+ \hat{a}  \right)  \,,  
\end{eqnarray}
%% 
que permite cancelar el exceso de ruido del OL, cuya intensidad es 
usualmente mucho mayor que la del emisor. Entonces aproximamos el 
OL por un estado coherente de amplitud $\beta$ y fase $\phi$, es decir 
$\hat{b} = \beta e^{i\phi}$, obteniendo 
%
\begin{eqnarray}
\hat{n}_{cd} &=& i \beta \left( \hat{a}^+ e^{i\phi} - \hat{a} e^{-i \phi}  \right) 
	= \beta e^{i \pi/2} \left( \hat{a}^+ e^{i\phi} - \hat{a} e^{-i \phi}  \right)
	\nonumber \\
&=& \beta \left[ \hat{a}^+ e^{i(\phi +\pi/2)} 
	+ \hat{a} e^{-i\pi} e^{-i( \phi -\pi/2)}  \right] 	\nonumber \\ 
&=& \beta \left[ \hat{a}^+ e^{i(\phi +\pi/2)} 
	+ \hat{a} e^{-i( \phi +\pi/2)}  \right] 	\nonumber \\ 	
&=& 2 \beta \hat{X}_{\phi +\pi/2} = i2\beta \hat{X}_{\phi} \,. 
\end{eqnarray}
%% 
$\hat{X}_{\phi}$ son los operadores de cuadratura (posici\'on y momentum adimensionales). Al normalizar $\hat{n}_{cd}$ quedan solamente las 
amplitudes de las cuadraturas. 

%%%%%%%%%%%%%%%%%%%%%%%%%%%%%%%%
\subsection{EES para DHB} 
La EES que simula la DHm esta dada por 
$i \hbar \ket{\dot{\bar{\Psi}}_c} = {\mcH}_W \ket{\bar{\Psi}_c}$ 
\cite{Carmichael93,WiMi93b}, donde 
%
\begin{eqnarray} 	\label{eq:HamiltonianHDm}
\mcH_W &=& \mcH + i \left(2 \gamma \langle \sigma_{\phi}  \rangle_c  
	+ \sqrt{\gamma} \,\eta_W(t) \right) e^{-i\phi}  \sigma_- \,,	
\end{eqnarray}
%% 
donde, para un \'atomo de dos niveles, 
%
\begin{eqnarray}
\mcH &=& H_0 - i \frac{\gamma}{2} \sigma_+ \sigma_-   \,, \quad 
	H_0 = \hbar \Delta \sigma_{ee} +\frac{1}{2} (\sigma_+ +\sigma_-) \,, 
	\nonumber \\ \\ 
\sigma_{\phi}  
	&=& \frac{1}{2}(\bar{\sigma}_+ e^{i\phi} +\bar{\sigma}_- e^{-i\phi} ) \,,
\end{eqnarray}
%% 
$\eta_W=dW/dt$ es un ruido blanco Gaussiano y $dW$ es un incremento 
de Wiener real y satisface 
%
\begin{equation} 	\label{eq:GaussianNoise}
\overline{dW} =0 \,, 	\quad \overline{dWdW} =dt \,, \quad 
	\overline{\eta_W(t) \eta_W(t')} = \delta(t-t') \,.
\end{equation}
%%

Podemos escribir la EES para DHm como 
%
\begin{subequations}
\begin{eqnarray}
\frac{d \ket{\bar{\Psi}_c}}{dt} 
	&=& \left\{ -i \mcH   +I_c^W(t) \sigma_- \right\}  \ket{\bar{\Psi}_c}   \,,	 \\
I_c^W(t) &=& \left(2 \gamma \langle \sigma_{\phi}  \rangle_c  
	+ \sqrt{\gamma} \,\eta_W(t) \right) e^{-i\phi} 		\\
&=& \gamma \left( \vxp{\bar{\sigma}_+}_c  
	+e^{-2i\phi} \vxp{\bar{\sigma}_-}_c \right) 
	+e^{-i\phi} \sqrt{\gamma} \, \eta_W(t)	\,,	\nonumber 
	\nonumber \label{HomodyneSSE} 
\end{eqnarray}
\end{subequations}
%%
donde $I_c^W$ es la fotocorriente medida. 

Vemos que la EES no es lineal debido a la presencia del valor esperado 
$\langle \sigma_{\phi} \rangle$ y que no contiene directamente 
informaci\'on de la intensidad y frecuencia del OL. Aparte de la fase 
$\phi$, el OL se presenta en la forma en que la estocasticidad entra: en 
lugar de evaluar si se registr\'o un fotoelectr\'on o no cada $\delta t$, 
ahora hay muchos fotoelectrones en dicho intervalo, produciendo un 
cambio cont\'inuo e infinitesimal en la funci\'on de onda. Entonces, m\'as 
que fotoelectrones aislados, se tiene una fotocorriente cuyas fluctuaciones 
tienen un car\'acter diferente al de fotoconteo: pasamos de ruido 
Poissoniano o sub-Poissoniano a ruido Gausiano. El procedimiento para 
llegar a la EES es muy largo y complicado. Esperamos poder reproducirlo 
en el futuro cercano. 

\textcolor{red}{Obtener: Correlaci\'on de  fotocorrientes} (seg\'un 
\cite{WiMi93b}): 
%
\begin{eqnarray} 	 \label{HomodyneAutoCorr}
\overline{ I_c(t+\tau) I_c(t) } &=& \gamma^2 
	\langle : \sigma_{\phi} (t+\tau) \sigma_{\phi} (t) : \rangle 
	+\gamma \delta (\tau) \,, \qquad
\end{eqnarray}
%%
donde usamos la f\'ormula cu\'antica de regresi\'on (??) y la 
Ec.(\ref{eq:GaussianNoise}). El espectro de fluorescencia se obtiene 
calculando la transformada de Fourier de esta correlaci\'on,
%
\begin{eqnarray} 	 \label{SpectrumHom}
S(\omega,\phi)-1 \propto \gamma \int_0^{\infty} dt\, \cos{\omega \tau} 
	\langle : \sigma_{\phi} (t+\tau) \sigma_{\phi} (t) : \rangle \,. \quad
\end{eqnarray}
%%
Este espectro es tambi\'en llamado espectro de \textit{squeezing}, 
el cual puede ser negativo, indicando la reducci\'on del ruido por debajo 
del nivel de ruido de disparo. 

%%%%%%%%%%%%%%%%%%%%%%%%%%%%%%%%
\subsection{Equations for 2LA resonance fluorescence} 
Las ecuacions para las amplitudes de probabilidad son  
%
\begin{subequations}
\label{eq:HomodyneSSEamp}
\begin{eqnarray}
dc_g &=& \left[ i\frac{\Delta}{2} c_g -i\frac{\Omega}{2}  
	c_e +\gamma \left( \vxp{\bar{\sigma}_-} e^{-2i\phi} 
	+\vxp{\bar{\sigma}_+} \right) c_e  \right] dt  	\nonumber \\
&& +\sqrt{\gamma} \, dW e^{-i\phi} c_e 	\\
%
dc_e &=& -\left[ i\frac{\Omega}{2} c_g 
 	+i\frac{\Delta}{2} c_e +\frac{\gamma}{2} c_e \right] dt
\end{eqnarray}
\end{subequations}
%%
La elecci\'on de $\phi=0$ (en fase) lleva a  
%
\begin{eqnarray*}
dc_g &=& \left[ i\frac{\Delta}{2} c_g -i\frac{\Omega}{2}  
	c_e +\gamma \vxp{\bar{\sigma}_X}  c_e \right] dt 
 	+\sqrt{\gamma} \, dW c_e  \,, \nonumber \\ 
 \end{eqnarray*}
%%
mientras que $\phi=\pi/2$ (en cuadrature) lleva a  
%
\begin{equation*}
dc_g = \left[ i\frac{\Delta}{2} c_g -i\frac{\Omega}{2} 
 	c_e -i\gamma \vxp{\bar{\sigma}_Y}  c_e \right] dt 
	 -\sqrt{\gamma} \, dW c_e 	\,.
\end{equation*}
%%
donde la nonlinealidad viene del valor de expectaci\'on de la 
dispersi\'on, $u= \vxp{\bar{\sigma}_X}$ y de la absorci\'on, 
$\vxp{\bar{\sigma}_Y}$, respectivamente. 

% (put a bar on Psi to indicate it is not normalized)

%%%%%%%%10%%%%%%%%20%%%%%%%%30%%%%%%%%40%%%%%%%%50%%%%%%%%60%%%%%%%%70%%%
\section{Detecci\'on Heterodina }
En la detecci\'on heterodina (DHt) la frecuencia del oscillador local es 
(mucho mayor? o) muy diferente que las frecuencias de los fotones 
emitidos y que el campo de excitaci\'on. A diferencia del caso de DHm, 
en lugar de un elemento \'optico que cambie la fase del OL, ahora se 
tiene un elemento que cambia su frecuencia. Entonces, el t\'ermino de 
variaci\'on r\'apida de la fase $e^{-2i\phi} \vxp{\sigma_-} \sigma_-$ en 
la Ec.(\ref{eq:HamiltonianHDm}) promedia a cero en un ciclo \'optico. 
Lo que sigue es un tratamiento muy simplificado de la Ref.\cite{WiMi93b}.
La EES que simula la DHt viene dada por 
$i \ket{\dot{\Psi}_c} ={\cal H}_Z \ket{\Psi_c}$, donde  
%
\begin{eqnarray} 	%\label{HeterodyneSSE} 
{\cal H}_Z &=& {\cal H} +i\left(\gamma \vxp{\bar{\sigma}_+}_c 
	+\sqrt{\gamma}  \eta_Z(t) \right) \sigma_- \,, 
\end{eqnarray}
%%
donde $\eta_Z=dZ/dt$ y  
%
\begin{eqnarray}
dZ &=& e^{i\omega^{\prime} t} dW =\frac{1}{\sqrt{2}} (dW_1 +idW_2)
\end{eqnarray}
%%
es un incremento de Wiener complejo, donde $dW_1, dW_2$ son 
incrementos de Wiener reales e independientes, que satisface
%
\begin{eqnarray} 	\label{eq:ComplexGaussNoise}
\overline{dZ} = 0 \,, \qquad \overline{dZdZ} &=&0  \,, 
	\qquad \overline{dZ^\ast dZ} =dt \,. 	\nonumber \\
	\overline{\eta_Z(t) \eta_Z(t')} &=& \delta(t-t') \,. 
\end{eqnarray}
%%
Alternativamente, la EES para detecci\'on heterodina se puede 
escribir como   
% 
\begin{subequations}
\begin{eqnarray}
\frac{d \ket{\Psi_c}}{dt} 
	&=& \left\{ -i \mcH   +I_c^Z(t) \sigma_- \right\}  \ket{\Psi_c}  \,, \\
I_c^Z(t) &=& \gamma  \vxp{\bar{\sigma}_+}_c   
	+\sqrt{\gamma} \, \eta_Z	\,,	\label{HeterodyneSSE}
\end{eqnarray}
\end{subequations}
%%
donde $I_c^Z$ es la fotocorriente medida. 

\textcolor{red}{NEXT: Correlaci\'on de  fotocorrientes} (seg\'un \cite{WiMi93b}): 
%
\begin{eqnarray} 	 \label{HomodyneAutoCorr}
\overline{ I_c(t+\tau) I_c(t) } &=& \gamma^2 
	\langle : \sigma_+ (t+\tau) \sigma_- (t) : \rangle 
	+\gamma \delta (\tau) \,, \qquad
\end{eqnarray}
%%
donde usamos la f\'ormula cu\'antica de regresi\'on (??) y la 
Ec.(\ref{eq:ComplexGaussNoise}). El espectro de fluorescencia se 
obtiene calculando la transformada de Fourier de esta correlaci\'on, 
%
\begin{eqnarray} 	 \label{SpectrumHom}
S(\omega)-1 \propto \gamma \int_0^{\infty} dt\, \cos{\omega \tau} 
	\langle  \sigma_+ (t+\tau) \sigma_- (t)  \rangle \,. \quad
\end{eqnarray}
%%

\subsection{Equations for 2LA resonance fluorescence} %%%%%%%%
The equations for the probability amplitudes are 
%
\begin{subequations}
\label{eq:HeterodyneSSEamp}
\begin{eqnarray}
dc_g &=& \left[  -i\frac{\Omega}{2} c_e  
	+\gamma \vxp{\bar{\sigma}_+} c_e \right] dt  
	+\sqrt{\gamma} \, dZ c_e \,, 	\quad \\
dc_e &=& -\left[ i\frac{\Omega}{2} c_g 
 	+i\frac{\Delta}{2} c_e +\frac{\gamma}{2} c_e \right] dt \,.
\end{eqnarray}
\end{subequations}
%%

%Use Carmichael's scheme of coupling the SSE with the equation for 
%the current, which includes the finite bandwidth of the detector in 
%the homodyne and heterodyne cases. 

%%%%%%%%10%%%%%%%%20%%%%%%%%30%%%%%%%%40
%%%%%%%%10%%%%%%%%20%%%%%%%%30%%%%%%%%40
\section{A Single SSE for QJ and QSD}
Carmichael, in Eq.(11) of \cite{Carmichael96}, has an SSE that includes 
jump, homodyne and heterodyne trajectories. It is more complicated, so 
we will try writing a simpler SSE: 
%
\begin{eqnarray} 	\label{GeneralSSE} 
d \ket{\bar{\Psi}_c } &=& \left\{ -i\mcH dt +\left[ 
	\gamma dt  \vxp{ \sigma_- e^{-i(\phi -\omega_{lo}t)} 
  	+\sigma_+ e^{i(\phi -\omega_{lo}t)} }_c \right. \right. 
		\nonumber 	\\
&& \left. \left. +\sqrt{\gamma} \, dW \right] 
	\sigma_- e^{-i(\phi-\omega_{lo}t)} 	\right\}   \ket{\bar{\Psi}_c} 	\,,	
\end{eqnarray}
%%
The local oscillator is characterized by the phase $\phi$, frequency 
$\omega_{lo}$, and a real Wiener increment $dW$. The properties of 
the noise term depend on the particular measurement scheme. For the 
heterodyne case, the average of phase terms is zero; for the homodyne 
case: $\omega_{lo} = \omega$. We go further, deleting the LO and 
nonlinear terms to obtain the SSE of quantum jumps. 

We now show that equation (citation?)
%
\begin{eqnarray}
d|\Psi \rangle &=& \left\{ \left[ \frac{\mathcal{H}}{i\hbar} 
	- \frac{\gamma}{2}(b^\dag b +r_+r_- -2r_+b) \right] \,dt 
	\right. \nonumber \\
	&& \left.  - \sqrt{\gamma}(r_- -b)e^{i\phi} dW \right\} |\Psi \rangle  \,,
\end{eqnarray}
%%
is a correct unraveling of the ME. Since the calculation does not require 
details of $\mathcal{H}$, we substitute the spin operators with boson 
operators $\sigma_+ =b^\dag, \, \sigma_- =b$. Further, we make 
$|\Psi_c\rangle =|\Psi \rangle$, $\langle \sigma_\pm \rangle_c =r_\pm$,
and $e^{i(\omega^\prime t +\theta)} =e^{i\phi}$. The master equation is 
recovered as \textcolor{red}{REVISAR}
%
\begin{widetext}
%
\begin{eqnarray*}
d(|\Psi\rangle \langle \Psi|) &=&
	(d|\Psi\rangle) \langle \Psi| +|\Psi\rangle (d\langle \Psi|) 
	+(d|\Psi\rangle) (d\langle \Psi|) 				\\
%
&=& \left[ \frac{H_S}{i\hbar} -\frac{\gamma}{2}(b^\dag b +r_+r_- -2r_+b) \right] 
	|\Psi \rangle \langle \Psi| \,dt  	
   -\sqrt{\gamma}(r_- -b)e^{i\phi} |\Psi \rangle\langle \Psi| \,dW \\
&& + |\Psi \rangle \langle \Psi| \left[ \frac{-H_S}{i\hbar} 
	-\frac{\gamma}{2}(b^\dag b +r_+r_- -2r_-b^\dag) \right] \,dt 
 -|\Psi \rangle\langle \Psi| \sqrt{\gamma}(r_+ -b^\dag)e^{-i\phi}\,dW \\
&&	+\left\{ \left[ \frac{H_S}{i\hbar} 
	-\frac{\gamma}{2}(b^\dag b +r_+r_- -2r_+b) \right] |\Psi \rangle \,dt 
	-\sqrt{\gamma}(r_- -b)e^{i\phi} |\Psi \rangle  dW	\right\}  \\
&&	\times \left\{ \langle \Psi| \left [\frac{-H_S}{i\hbar} 				
	-\frac{\gamma}{2}(b^\dag b +r_+r_- -2r_-b^\dag) \right] \,dt 
  -\langle\Psi| \sqrt{\gamma}(r_+ -b^\dag)e^{-i\phi} \,dW \right\} \\
%
&=& \frac{1}{i\hbar} \left[ H_S, |\Psi\rangle \langle \Psi| \right] \,dt 
	-\frac{\gamma}{2}\,dt \left\{  b^\dag b |\Psi\rangle \langle \Psi| 
	+|\Psi \rangle \langle \Psi| b^\dag b 
  +2r_+r_- |\Psi\rangle\langle\Psi| -2r_+b |\Psi\rangle\langle\Psi| 
	-2r_- |\Psi\rangle\langle\Psi| b^\dag \right\}		\\
&&	+\gamma dW^2 (r_- -b)|\Psi\rangle\langle\Psi|(r_+ -b^\dag)
	+O(dt^2) +O(dW) \,.
\end{eqnarray*}
%%
\end{widetext}
Now, neglecting terms $\sim dt^2$, performing the ensemble average 
on $\rho= \overline{|\Psi\rangle\langle\Psi|}$ with the (Ito?) rules 
$\overline{dW}=0$ and $\overline{dW^2}=dt$, and with terms that 
cancel each other we obtain the ME 
%
\begin{equation} 	\label{ME-RF}
\dot{\rho}(t) = \frac{1}{i\hbar} [H_S, \rho] 
+\frac{\gamma}{2} ( 2b\rho b^\dag -b^\dag b \rho -\rho b^\dag b ) \,.		
\end{equation}
%%

%%%%%%%%%%%%%%%%%%%%%%%%%%%%%%%%%
\subsection{Equations for 2LA resonance fluorescence} 
We find the equations for the (unnormalized) probability amplitudes 
for the ground and excited states in the frame of free oscillations 
%
\begin{subequations}
\label{eq:GeneralSSEamp}
\begin{eqnarray}
dc_g &=& +\sqrt{\gamma} \, e^{-i(\phi-\Delta_{lo}t)} dW c_e 
	+\left[ i\frac{\Delta}{2} c_g -i\frac{\Omega}{2} c_e 	\right. 	
	\nonumber 	\\
&& \left. +\gamma \left( \vxp{\bar{\sigma}_+} 
	+\vxp{\bar{\sigma}_-} e^{-2i(\phi-\Delta_{lo}t)} \right) c_e \right] dt 	\\
%
dc_e &=& -\left[ i\frac{\Omega}{2} c_g 
 	+i\frac{\Delta}{2} c_e +\frac{\gamma}{2} c_e \right] dt
\end{eqnarray}
\end{subequations}
%%


%%%%%%%%10%%%%%%%%20%%%%%%%%30%%%%%%%%40%%%%%%%%50%%%%%%%%60%%%%%%%%70%%%
\section{Conditional Homodyne Detection}
The derivation of the SSE is very long hence we show a few steps only. Within the framework of quantum trajectory theory CHD is described by a quantum stochastic process comprised of a measurement record, and the quantum state of the source, 
$\ket{\psi_{\rm R}}$, conditioned on that record. An ensemble of realizations, 
$\ket{\psi_{{\rm R}_{(n)}}^{(n)} }$, $n=1,2,\cdots$, provide an unraveling of the 
source density operator; the associated records are realizations of laboratory data \cite{CCFO00,FOCC00}. We analyze the data record for a single, continuous 
realization, and construct the BHD current average from its conditional sampling 
over time.

The homodyne current is the stochastic process
%
\begin{equation} 	\label{eqn:current} 
dI = - \Gamma(Idt-dQ_t),
\end{equation}
%%
with
%
\begin{equation}
dQ_t = 2 \sqrt{2r} \vxp{b_\theta}_{\rm R} dt +dW_t,
\end{equation}
%% 
and
%
\begin{equation}
 b_\theta \equiv ( b e^{-i\theta} +{\rm H.c.})/2, 	\qquad
	 b\equiv a+Ae^{i\phi},
\end{equation}
%%
where $a$ is the source-field annihilation operator, $Ae^{i\phi}$ 
is the complex amplitude of the coherent offset, $r$ is the reflection
coefficient at the correlator input (BS2), $\theta$ is the local 
oscillator (LO) phase, and $\Gamma$ is the BHD bandwidth; $dW_t$ is a 
Wiener process modeling infinite bandwidth shot noise. The units of 
time are inverse bandwidths (half-widths) of the quasimonochromatic 
source.

The conditional state of the source (unnormalized) obeys
%
\begin{eqnarray}
d \ket{\bar\psi}_{\rm R} &=& \left[-(i\mcH +2Ae^{-i\phi} a)dt 
	+\sqrt{2r} b e^{-i\theta}dQ_t \right] \ket{\bar\psi}_{\rm R} \,, 
	\nonumber \\ 	\label{eqn:schrodinger} 
\end{eqnarray}
%%
where $\mcH$ is the non-Hermitian Hamiltonian and the other terms 
account for the measurement backaction. There are also quantum jumps,
%
\begin{equation}
\ket{\bar\psi}_{\rm R} \to b \ket{\bar\psi}_{\rm R},
\end{equation}
%% 
at every firing of the ``start'' detector, where the instantaneous
firing rate is
%
\begin{equation}
R_s = 2(1-r) \vxp{n_b}_{\rm R}, 	\qquad n_b\equiv b^\dagger b.
\end{equation}
%%

We integrate Eq.\ (\ref{eqn:current}) formally and average $N_s$ 
samples of the photocurrent with ``starts'' at times $t_j$, 
$j=1,\ldots,N_s$, to obtain the raw signal
%
\begin{equation}
H(\tau) = \frac\Gamma{N_s}\sum_{j=1}^{N_s}\int_{-\infty}^{t_j+\tau}\!dQ_t\,
	e^{\Gamma(t-t_j-\tau)}.
\end{equation}
%%
For large $N_s$, the sum over $j$ amounts to an average over past and
future trajectory records, evaluate as \cite{CCFO00}
%
%\begin{mathletters}
\begin{equation}
\frac{H(\tau)}{2\sqrt{2r}} =
  \Gamma \int_{-\infty}^\tau dt\,\frac{ \vxp{ : n_b(0) b_\theta(\tau) :}}{\vxp{n_b}} 			e^{\Gamma(t-\tau)} +\Xi(\tau),
\label{eqn:largecor} 
\end{equation}
%%
%
\begin{equation}
\overline{\Xi(\tau)\Xi(\tau^\prime)} 
	= N_s^{-1}\Gamma e^{-\Gamma |\tau-\tau^\prime|},
\end{equation}
%%
where $::$ denotes time and normal ordering and $\Xi(\tau)$ is the
residual shot noise.
%\end{mathletters}
%%
From Eq.(\ref{eqn:largecor}), CHD measures a third-order correlation
function of the electromagnetic field, cross-correlating 
photodetections with quadrature amplitude fluctuations.

%%%%%%%%10%%%%%%%%20%%%%%%%%30%%%%%%%%40%%%%%%%%50%%%%%%%%60%%%%%%%%70%%%
\section{Quadrature Fluctuations and Spectrum of Squeezing in 
Cavity QED}
The squeezing connection is made by separating the mean field 
amplitude, $\beta$, from the source-field fluctuation, $\Delta a$, 
$\langle \Delta a \rangle =0$. We write
%
\begin{equation}
 b=\beta+\Delta a, 		\qquad 
\beta\equiv\langle b\rangle = \langle a \rangle +Ae^{i\phi}.
\end{equation}
%%
%
\begin{eqnarray}
b_\theta &=& \beta_\theta +\Delta q_\theta \,,	\nonumber	\\
\beta_\theta &=& \langle b \rangle = 
	(\beta e^{-i\theta} +\beta^\ast e^{i\theta} )/2 \,, \nonumber	\\
\Delta q_\theta &=& 
  (\Delta a e^{-i\theta} +\Delta a^\dag e^{i\theta} )/2 \,.
\end{eqnarray}
%%
The photon number at the particle detector is 
%
\begin{eqnarray}
n_b = b^\dag b &=& |\beta|^2 +\Delta a^\dag\Delta a  
	+\beta^\ast\Delta a +\beta\Delta a^\dag	\,,  \nonumber		\\
\langle b^\dag b\rangle &=& 
	|\beta|^2 + \langle \Delta a^\dag \Delta a \rangle \,.
\end{eqnarray}
%%
We assume third-order moments in $\Delta a^\dagger$ and $\Delta a$ 
vanish: the fluctuations might be Gaussian, or the assumption is
approximately valid because the source field is weak. Therefore, 
%
\begin{eqnarray}
\lefteqn{
\frac{\langle\,:\! n_b(0) b_\theta(\tau)\!:\,\rangle}{\langle n_b\rangle} 
	= \beta_\theta } \hspace{1cm} \nonumber\\ 
	&& +\frac{\langle\,:\! [\beta^\ast \Delta a(0) 
	+\beta\Delta a^\dag(0)] \Delta q_\theta(\tau)\!:\,\rangle}
	{\langle n_b\rangle} +\xi(\tau),		\nonumber\\ 				
\end{eqnarray}
%%
Noting that 
$\langle\,:\! \Delta a(0)\Delta a(\tau)\!:\,\rangle$ and 
$\langle\,:\! \Delta a^\dag(0)\Delta a^\dag(\tau)\!:\,\rangle$ 
only contribute noise, and that (?) 
$\langle\,:\! \Delta a(0)\Delta a^\dag(\tau)\!:\,\rangle = 
\langle\,:\! \Delta a^\dag(\tau)\Delta a(0)\!:\,\rangle$  
we can write
%\begin{widetext}
%
\begin{eqnarray}
\frac{\langle\,:\! n_b(0) b_\theta(\tau)\!:\,\rangle}{\langle n_b\rangle} 
	&=& \xi(\tau) + \beta_{\theta} \nonumber \\
	&+& \frac{\beta^\ast e^{i\theta} +\beta e^{-i\theta} }{2} 
	 \frac{\langle\,:\! \Delta a^\dag (0) \Delta a (\tau)\!:\,\rangle}
	{\langle n_b\rangle}  \,, \nonumber \\	
\end{eqnarray}
%%
%\end{widetext}

Let us analyze two cases: (a) In cavity QED 
$\langle a \rangle \neq 0$, so the offset is not needed $A=0$, 
$\langle n_b \rangle = \langle n_a \rangle = 
|\alpha|^2 +\langle \Delta a^\dag \Delta a \rangle$, and 
$\beta=\beta^\ast=\alpha$, 
%
\begin{eqnarray}
\frac{\langle\,:\! n_a(0) a_\theta(\tau)\!:\,\rangle}
	{\langle n_a\rangle} &=& \alpha \cos{\theta}  \\
	&&+2\alpha \frac{\langle\,:\! \Delta q_{0^\circ}(0) 
	\Delta q_\theta(\tau)\!:\,\rangle}
	{\alpha^2 +\langle \Delta a^\dag \Delta a \rangle} 
	+\xi(\tau) \,,	\nonumber 
\end{eqnarray}
%%
%
\begin{eqnarray}
\frac{\langle\,:\! n_a(0) a_\theta(\tau)\!:\,\rangle}
	{\langle n_a\rangle} = \alpha \cos{\theta}  \left[ 1+  
\frac{\langle\,:\! \Delta a^\dag (0) \Delta a (\tau)\!:\,\rangle}
	{\alpha +\langle \Delta a^\dag \Delta a \rangle} \right]
 	+\xi(\tau) \,.			\nonumber
\end{eqnarray}
%%
(b) Lock the offset's phase to the LO's ($\phi=\theta$), so that 
$\beta_\theta = A+ \alpha \cos{\theta}$, and match $|\beta|$ 
to the R.M.S. fluctuation:
%
\begin{eqnarray}
\frac{\langle\,:\! n_b(0) b_\theta(\tau)\!:\,\rangle}
	{\langle n_b\rangle} &=& (A+ \alpha\cos{\theta}) \left[ 1+  
\frac{\langle\,:\! \Delta a^\dag (0) \Delta a (\tau)\!:\,\rangle}
	{|\beta|^2 +\langle \Delta a^\dag \Delta a \rangle} \right] 
	\nonumber \\
	&&+\xi(\tau),	 
\end{eqnarray}
%%

Thus, we set
%
\begin{equation}
\beta=\sqrt{\langle\Delta a^\dagger\Delta a\rangle}\,e^{i\theta}.
\end{equation}
%%
Then, for broadband detection ($\Gamma\gg1$), the normalized signal
$h_\theta(\tau)\equiv H(\tau)/2\sqrt{2r}|\beta|$ is the correlation 
function
%
%\begin{mathletters}
\begin{equation} 	\label{eqn:smallcor}
h_\theta(\tau) = 1+\frac{\langle\,:\Delta q_\theta(0)\Delta
	q_\theta(\tau)\!:\,\rangle}{\langle\Delta a^\dagger\Delta a\rangle}
	+\xi(\tau),
\end{equation}
%%
%
\begin{equation} 	\label{eqn:noise}
\overline{\xi(\tau)\xi(\tau^\prime)} = (16rN_s|\beta|^2)^{-1}
	\Gamma e^{-\Gamma|\tau-\tau^\prime|},
\end{equation}
%\end{mathletters}
%%
and CHD measures the autocorrelation of the quadrature amplitude
fluctuations. The squeezing spectrum \cite{CCFO00,FOCC00},
%
\begin{equation}
S_\theta(\omega)=4F\!\int_0^\infty\! 
	d\tau\cos(\omega\tau)[\overline{h_\theta (\tau)}-1],
\end{equation}
%%
where $F=2\kappa \langle n_b\rangle=4|\beta|^2$ is the photon flux into
the correlator and $\overline{h_\theta(\tau)}\equiv\lim_{N_s\to\infty}
h_\theta(\tau)$, is independent of the fraction of light,
$r$, sent to the BHD. Of course, for a fixed measurement time, the
signal-to-noise ratio does depend on $r$, which enters explicitly and
through $N_s$ into Eq.\ (\ref{eqn:noise}). [A separate efficiency 
enters the measurement of $F$.]

The standard signature of squeezing is a reduced shot noise level which
is interpreted as a deamplification of the vacuum fluctuations. CHD
provides another view. The focus is on the fluctuations of the light
emitted by the source. The shot noise is irrelevant, ideally being
eliminated through the conditional average. The fluctuations are
nonclassical. They violate inequalities that any classical field must
satisfy---$\langle\,:\!\Delta q_\theta(0)\Delta q_\theta(\tau)\!:
\,\rangle$ is classically incompatible with
%
\begin{equation} 	\label{eqn:variance}
\langle\Delta a^\dagger\Delta a\rangle = 
	\langle\,:\!(\Delta q_\theta)^2\!:\,\rangle+\langle\,:\!(\Delta
	q_{\theta+\pi/2})^2\!:\,\rangle.
\end{equation}
%%
Classically, the quadrature variances are positive. It follows,
trivially, from Eqs.\ (\ref{eqn:smallcor}) and (\ref{eqn:variance}), 
that
%
%\begin{mathletters}
\begin{equation} 	\label{eqn:inequality1}
0\leq\overline{h_\theta(0)} -1 \leq 1.
\end{equation}
%%
More generally one proves the Schwarz inequality
%
\begin{equation} \label{eqn:inequality2}
|\overline{h_\theta(\tau)}-1|\leq|\overline{h_\theta(0)}-1|\leq1.
\end{equation}
%\end{mathletters}
%%
One quadrature variance is negative, however, if there is squeezing 
below the quantum limit. It follows that the squeezed (unsqueezed)
fluctuations violate the lower (upper) bound of (\ref{eqn:inequality1}),
and fluctuations of either quadrature amplitude may violate
(\ref{eqn:inequality2}). Conditions of low photon flux are of 
particular interest. 

%%%%%%%%10%%%%%%%%20%%%%%%%%30%%%%%%%%40%%%%%%%%50%%%%%%%%60%%%%%%%%70%%%%%
\section{Quantum Trajectories for Spectra}
Unraveling of the master equation in the frequency domain, specifically 
the Tian-Carmichael \cite{TiCa92} and Holland \cite{Holland98} 
approaches. These methods calculate the number of photons emitted 
at a certain frequency and time. Methods based on two-time field 
correlations will probably be out of consideration (Dalibard et al, 
Zoller et al, etc). 

In the Tian-Carmichael method the light from the source is coupled, 
with strength $\epsilon \ll 1$, to the filter with zero- and one-photon 
states, $|0,1 \rangle_F$, half-width $\Gamma$, and resonance frequency 
$\nu$. The non-Hermitian Hamiltonian is 
%
\begin{eqnarray}
\mathcal{H} = H_{\mathrm{eff}} +(\nu -i\Gamma)b^{\dag}b 
	-i\sqrt{\epsilon \Gamma} b^{\dag} \sigma_-
\end{eqnarray}
%%
where $H_{\mathrm{eff}}$ is the effective Hamiltonian, $b^{\dag}$ and 
$b$ act only on the filter states. The unnormalized wave function is 
%
\begin{eqnarray}
|\overline{\Psi}_c (\nu,t)\rangle  &=& 
|\overline{\Psi}_c^0 (\nu,t) \rangle |0 \rangle_F 
+|\overline{\Psi}_c^1 (\nu,t)\rangle |1 \rangle_F \,, \qquad
\end{eqnarray}
%%
The Schr\"odinger equation splits into two sets of equations
%
\begin{subequations}
\begin{eqnarray}
|\dot{\overline{\Psi}_c^0} (\nu,t) \rangle 
	&=& - i H_{\mathrm{eff}} |\overline{\Psi}_c^0(\nu,t) \rangle \,, 	\\
%
|\dot{\overline{\Psi}_c^1}(\nu,t) \rangle &=& \left[ -i H_{\mathrm{eff}} 
 -(\Gamma +i\nu) \right] |\overline{\Psi}_c^1 (\nu,t) \rangle  \nonumber \\
&& -\sqrt{\epsilon \Gamma} \sigma_- |\overline{\Psi}_c^0 (\nu,t) 
	\rangle \,, 	
\end{eqnarray}
\end{subequations}
%%
The spectrum is given by 
%
\begin{eqnarray}
S(\nu,t) &=& 
\frac{ \langle \overline{\Psi}_c^1(\nu,t) |\overline{\Psi}_c^1(\nu,t) 
\rangle }{ \langle \overline{\Psi}_c(\nu,t) |\overline{\Psi}_c(\nu,t) 
\rangle } 
\end{eqnarray}
%%
which is the mean photon number in the filter cavity for different 
settings of the filter frequency $\nu$. This spectrum is shown to be the 
same as the EW spectrum. 

What is $|\overline{\Psi}_c^1 (\nu,t)\rangle$ like for a two-level atom?



%%%%%%%%10%%%%%%%%20%%%%%%%%30%%%%%%%%40%%%%%%%%50%%%%%%%%60%%%%%%%%70%%%
% \section{Conclusions}

%%%%%%%%10%%%%%%%%20%%%%%%%%30%%%%%%%%40%%%%%%%%50%%%%%%%%60%%%%%%%%70%%%
%\section*{Appendix}

%%%%%%%%10%%%%%%%%20    REFERENCES    40%%%%%%%%50%%%%%%%%60%%%%%%%%70%%%
\begin{references} 
%
\bibitem{Carmichael93} H. J. Carmichael \textit{An Open Systems
 Approach to Quantum Optics} (Springer, Berlin, 1993).
%
\bibitem{Carmichael02} H. J. Carmichael \textit{Statistical Methods in 
Quantum Optics 1: Master Equations and Fokker-Planck Equations} 
(Springer, Berlin, 1999). 
%
\bibitem{Carmichael07} H. J. Carmichael \textit{Statistical Methods in 
Quantum Optics 2: Non-Classical Fields} (Springer, Berlin, 2008). 
%
\bibitem{Carmichael96} H. J. Carmichael, Stochastic Schr\"odinger 
equations: What they mean and what they can do, in \textit{Coherence 
and Quantum  Optics}, J. H. Eberly, L. Mandel, and E. Wolf, eds. (Plenum 
Press, New York, 1996).
%
\bibitem{DaCM92} J. Dalibard, Y. Castin, and K. Molmer, Phys. Rev. 
Lett. \textbf{68}, 580 (1992).
%
\bibitem{MoCD92} K. Molmer, Y. Castin, and J. Dalibard, J. Opt. Soc. 
Am. B \textbf{10}, 524 (1993). 
%
\bibitem{DuZR92} R. Dum, P. Zoller, and H. Ritsch, Phys. Rev. A 
\textbf{45}, 4879 (1992).
%
\bibitem{PlKn98} M. B. Plenio and P.L. Knight, Rev. Mod. Phys. 
\textbf{70}, 11 (1998).
%
\bibitem{Percival98} I. Percival, \textit{Quantum State Diffusion}  
(Cambridge University Press, Cambridge, 1998). 
%
\bibitem{WiMi93a} H. M. Wiseman and G. J. Milburn, Phys. Rev. A 
\textbf{47}, 988 (1993).
%
\bibitem{WiMi93b} H. M. Wiseman and G. J. Milburn, Phys. Rev. A 
\textbf{47}, 1652 (1993).
%
\bibitem{WiMi11} H. M. Wiseman and G. J. Milburn, \textit{Quantum 
Measurement and Control} (Cambridge University Press, Cambridge, 
2011). 
%
\bibitem{GaZoQN} C. Gardiner and P. Zoller, \textit{Quantum Noise, 
3rd ed.} (Springer-Verlag, Berlin, 200?). 
%
\bibitem{GaZoQO1} C. Gardiner and P. Zoller, \textit{The Quantum World of 
Ultra-Cold Atoms and Light Book I: Foundations of Quantum Optics} 
(World Scientific, Singapore, 2014). 
%
\bibitem{GaZoQO2} C. Gardiner and P. Zoller, \textit{The Quantum World 
of Ultra-Cold Atoms and Light Book II: The Physics of Quantum-Optical 
Devices} (World Scientific, Singapore, 2015).  
%
\bibitem{BrPe02} H.-P. Breuer and F. Petruccione, \textit{The Theory of 
Open Quantum Systems} (Oxford University Press, Oxford, 2002). 
%
\bibitem{CCFO00} H. J. Carmichael, H. M. Castro-Beltran, G. T. Foster, 
and L. A. Orozco, Phys. Rev. Lett. \textbf{85}, 1855 (2000).
%
\bibitem{FOCC00} G. T. Foster, L. A. Orozco, H. M. Castro-Beltran, and 
	H. J. Carmichael, PRL \textbf{85}, 3149 (2000).
%
\bibitem{BeMa11} P. R. Berman and V. S. Malinovsky, \textit{Principles 
of Laser Spectroscopy and Quantum Optics} (Princeton University Press, 
Princeton, 2011).
%
\bibitem{Huard16} P. Campagne-Ibarcq, P. Six, L. Bretheau, A. Sarlette, 
M. Mirrahimi, P. Rouchon, and B. Huard, Observing quantum state 
diffusion by heterodyne detection of fluorescence, Phys. Rev. X 
\textbf{6}, 011002 (2016).
%
\bibitem{Carmichael93casc} H. J. Carmichael, Quantum trajectory 
theory for cascaded open systems,  Phys. Rev. Lett. \textbf{70}, 2273 
(1993). 
%%
\bibitem{TiCa92} L. Tian and H. J. Carmichael, Phys. Rev. A \textbf{46}, 
	R6801 (1992).
%
\bibitem{Holland98} M.~J. Holland, Phys. Rev. Lett. \textbf{81}, 5117 
(1998); Comment by D.~B. Horoshko and S.~Ya. Kilin, Phys. Rev. Lett. 
\textbf{83}, 2469 (1999); 
Reply by M.~J. Holland, Phys. Rev. Lett. \textbf{83}, 2470 (1999).
%
\bibitem{CaGH15} H. M. Castro-Beltran, L. Guti\'errez, and L. Horvath,  
Appl. Math. Inf. Sci. Opt. \textbf{9}, 2849 (2015). 
%
\bibitem{KiKn89} M. S. Kim and P. L. Knight, Phys. Rev. A 
\textbf{40}, 215 (1989).
%
\bibitem{CaWa76} H. J.  Carmichael and D.F. Walls, J. Phys. B 
	{\bf 9}, 1199 (1976). 
%
\bibitem{KiMa76} H.J. Kimble and L. Mandel, Phys. Rev. A 
\textbf{13}, 2123 (1976).
%
\bibitem{antibunch} H. J. Kimble, M. Dagenais, and L. Mandel, 
	Phys. Rev. Lett. {\bf 39}, 691 (1977).
%F. Diedrich and H. Walther, {\it ibid} {\bf 58}, 203 (1987); 
%V. Gomer, F. Strauch, B. Ueberholz, S. Knappe, and D. Meschede, 
%	Phys. Rev. A {\bf 58}, R1657 (1998).
%
\bibitem{CSVR89} H. J. Carmichael, S. Singh, R. Vyas, and P. R. Rice, 
Photoelectron waiting times and atomic state reduction in resonance 
fluorescence, Phys. Rev. A \textbf{39}, 1200 (1989).  

 
\end{references}


\end{document}

%%%%%%%%%%%%%%%%%%%%%%%%%%%%%%%%%%%%
%%%%%%%%%%%%%%%%%%%%%%%%%%%%%%%%%%%%


